\documentclass[11pt]{article}
\usepackage[a4paper,margin=27mm]{geometry}
\usepackage[T1]{fontenc}
\usepackage[utf8]{inputenc}
\usepackage{lmodern,amsmath,amssymb,amsthm,mathtools,booktabs,longtable,array}
\usepackage[colorlinks=true,linkcolor=blue,citecolor=blue,urlcolor=blue]{hyperref}
\usepackage{microtype}
\setlength{\emergencystretch}{2em}
\newtheorem{theorem}{Theorem}[section]
\newtheorem{proposition}[theorem]{Proposition}
\newtheorem{lemma}[theorem]{Lemma}
\newtheorem{corollary}[theorem]{Corollary}
\theoremstyle{definition}\newtheorem{definition}[theorem]{Definition}
\theoremstyle{remark}\newtheorem{remark}[theorem]{Remark}
\newcommand{\E}{\mathbb E}
\newcommand{\Prb}{\mathbb P}
\newcommand{\Z}{\mathbb Z}
\newcommand{\rad}{\operatorname{rad}}
\newcommand{\unit}{\mathbf1}
\newcommand{\floor}[1]{\left\lfloor #1\right\rfloor}
\newcommand{\ceil}[1]{\left\lceil #1\right\rceil}
\newcommand{\gcds}{\operatorname{gcd}}
\newcommand{\Li}{\operatorname{Li}_2}
\DeclareMathOperator{\supp}{supp}
\title{Mean rational complexity in integer bases:\\
a uniform bound and a two-prime theorem}
\author{Oleksiy Babanskyy}
\hypersetup{
  pdftitle={Mean rational complexity in integer bases: a uniform bound and a two-prime theorem},
  pdfauthor={Oleksiy Babanskyy}
}
\date{}
\begin{document}
\maketitle
\begin{abstract}
For a uniformly chosen residue modulo \(B^N\), let its rational height be
the smallest \(\max(|a|,d)\) among congruent rational representations
\(a\equiv sd\pmod {B^N}\), with \(d>0\) coprime to the entire modulus.
We prove a uniform upper bound strictly below \(N/2+49/50\) for the mean
base-\(B\) logarithm of this height. The additive correction can be removed
for prime-power bases, and is strictly negative for bases with exactly two
distinct prime factors. The latter result combines explicit estimates with
finite integer certificates covering arbitrary positive valuations.
The correction cannot be removed for all bases: \(B=30,N=2\) is a certified
counterexample. An explicit lower estimate also gives an eventually
positive family with radical \(510510\).
The self-contained argument places the optimal global additive correction
strictly between zero and \(49/50\); a classical Jacobsthal bound further
gives a correction strictly below \(4/5\).
Its supremum reduces to finitely many bases, though neither its value
nor finite-length attainment is determined here.
Classical rational-reconstruction and lattice tools are distinguished
from the new bounds. The companion specifies retained evidence,
independent checks and full reconstruction commands.
\end{abstract}

\section{Model, results and antecedents}
Throughout, logarithms without a subscript are natural logarithms.
All probability spaces in this paper contain \emph{all} residues, not only
invertible states.

\begin{definition}
For \(M\ge2\) and \(s\in\Z/M\Z\), set
\[
 R_M(s)=\min_{\substack{d\in\Z_{>0},\ \gcd(d,M)=1\\a\in\Z,\ a\equiv sd\ (M)}}
          \max(|a|,d).
\]
For integers \(B\ge2,N\ge1\), write \(R_{B,N}=R_{B^N}\), with \(s\)
uniformly distributed on \(\Z/B^N\Z\), and put
\[
 E_{B,N}=\E\log_B R_{B,N},\qquad
 \mu(M)=\E\log(R_M/\sqrt M).
\]
Thus \(E_{B,N}-N/2=\mu(B^N)/\log B\).
\end{definition}
The minimum exists: \(d=1\) with a centered representative of \(s\) gives
\(R_M(s)\le\max(1,\floor{M/2})\). In particular \(R_M(0)=1\).
An admissible witness can be made primitive by dividing \(a,d\) by their
gcd; the divisor is coprime to \(M\). The full-unit condition on \(d\) is
essential.

\begin{theorem}\label{thm:main}
For every \(N\ge1\):
\begin{enumerate}
\item For every integer \(B\ge2\), \(E_{B,N}<N/2+49/50\).
\item If \(B\) is a power of a prime, \(E_{B,N}\le N/2\).
\item If \(B\) has exactly two distinct prime factors, \(E_{B,N}<N/2\).
\end{enumerate}
\end{theorem}
The binary case of the second assertion is the theorem of
\cite{binary}; its proof is imported and attributed here. The odd-prime
argument below is independent of limiting laws. The two-prime theorem
includes \(B=p^a r^b\) with arbitrary \(a,b\ge1\); it is not restricted to
squarefree bases or equal valuations.

\begin{theorem}\label{thm:positive}
One has \(E_{30,2}>1\). If \(\rad(B)=q_0=510510\) and \(B^N\ge q_0^2\),
then
\[
 E_{B,N}>N/2+\frac{1}{1000\log B}.
\]
Consequently the zero-correction bound does not hold for all bases.
\end{theorem}

\begin{theorem}\label{thm:optimal}
Set
\[
 c(B)=\sup_{N\ge1}(E_{B,N}-N/2),\quad c_+(B)=\max(0,c(B)),\quad
 C^*=\sup_{B\ge2}c(B).
\]
Then \(c_+(B)\to0\) as \(B\to\infty\), and \(0<C^*<49/50\).
There exists a finite \(B_0\) such that
\(C^*=\max_{2\le B\le B_0}c(B)\).
This does not assert that the supremum over \(N\) is attained.
\end{theorem}

Question 3.2 of the survey \cite[p.~20282]{survey} asks for an upper bound
in general bases. Our answer concerns the explicitly defined rational
height model. A connection integer in an FCSR realization may have the
additional condition \(d\equiv-1\pmod B\), and register lengths may be
rounded. We do not identify such quantities with this minimum for
composite bases without a further equivalence theorem.
The cited author draft of the monograph \cite[Sections 19.3,21.3,21.5]{GK}
treats coprime rational reconstruction and complexity estimates; its
Section 19.1 uses a different connection-integer convention.
Section locators here refer to that author draft, not a verified
pagination of the published edition. The known lower estimates already
give the leading coefficient \(1/2\) for fixed bases. We do not claim that
coefficient as a new result.

Chen--Winterhof \cite{CW} gives probabilistic binary results; its accessible
preprint version has been checked, while the final published text has not
been fully accessed for the source comparison reported here. The exact binary mean and
variance results in \cite{binary} are prior results, not new claims here.
Farey fractions with excluded denominators were studied by
Boca--Heersink--Spiegelhalter \cite{BHS}; congruence sections by
Heersink \cite{Heersink}; smallest-denominator laws and logarithmic
moments by Marklof \cite{Marklof,MarklofLog}. Those results are relevant
antecedents, but an unmarked law cannot simply be rescaled to impose all
the full-unit restrictions. Priority of the present bounds is not certified
by the source comparison presented here.

\section{A lattice upper bound}
For a squarefree integer \(q\), let \(J(q)\) be the maximum gap between
successive integers coprime to \(q\), including the cyclic gap modulo \(q\).
Let \(\lambda_s\) be the shortest Euclidean length of a nonzero vector in
\[
 \Lambda_s=\{(a,d)\in\Z^2:a\equiv sd\pmod M\},
 \qquad \det\Lambda_s=M.
\]

\begin{lemma}\label{lem:wheel}
For \(q=\rad(M)\),
\[
 R_M(s)\le\sqrt{1+J(q)^2/3}\,\frac M{\lambda_s}.
\]
Also \(J(q)\le2\cdot4^{\omega(q)}\), uniformly over squarefree \(q\).
\end{lemma}
\begin{proof}
A shortest vector \(u\) is primitive in the lattice and extends to a
basis \(u,v\). After subtracting an integer multiple of \(u\), write
\(v=\theta u+v_\perp\), where \(|\theta|\le1/2\) and
\(\|v_\perp\|=M/\lambda_s\). Minimality of \(u\) gives
\(\lambda_s^4\le4M^2/3\).
The second coordinates of \(u,v\) have gcd one, since the lattice contains
\((s,1)\). If a prime dividing \(q\) divides the second coordinate of \(u\),
it cannot divide that of \(v+ju\). For every other prime, precisely one
residue class of \(j\) is forbidden. By CRT this is a translate of the
unit pattern modulo a divisor of \(q\). Its maximum gap is at most \(J(q)\).
Choose an allowed integer \(j\) within \(J(q)/2\) of \(-\theta\).
Then \(v+ju\) has a full-unit second coordinate and
\[
 \|v+ju\|^2\le M^2/\lambda_s^2+(J(q)^2/4)\lambda_s^2
 \le (1+J(q)^2/3)M^2/\lambda_s^2.
\]
Changing the sign makes the denominator positive.
If the denominator of \(u\) is already a unit, \(u\) itself is available;
the same bound holds because
\(\lambda_s^2\le2M/\sqrt3\) and \(J(q)\ge1\).

For the second assertion, inclusion--exclusion shows that an interval of
\(L\) consecutive integers contains at least
\(L\varphi(q)/q-2^{\omega(q)}\) units. Since
\(\varphi(q)/q\ge2^{-\omega(q)}\), \(L=2\cdot4^{\omega(q)}\)
makes this quantity positive. Every such interval contains a unit.
\end{proof}

\begin{lemma}\label{lem:shorttail}
For \(1\le Y<M\),
\(\Prb(\lambda_s\le Y)\le3Y^2/M\).
Consequently, with \(Z_M=R_M/\sqrt M\),
\[
 \Prb(Z_M>z)\le\min(1,(J(q)^2+3)z^{-2}),\qquad
 \mu(M)\le \frac{\log(J(q)^2+3)+1}{2}.                 \tag{2.1}\label{eq:wheelmean}
\]
\end{lemma}
\begin{proof}
A horizontal nonzero lattice vector has length at least \(M\).
For \(1\le d\le Y\), put \(g=\gcd(d,M)\).
An integer \(a\) is incident with \(g\) residues \(s\) if \(g\mid a\),
and with none otherwise. Hence the number of incidences in the disk is
at most
\[
 \sum_{1\le d\le Y}\bigl(2\sqrt{Y^2-d^2}+\gcd(d,M)\bigr)
 \le (1+\pi/2)Y^2<3Y^2.
\]
Here \(\gcd(d,M)\le d\), \(\sum_{d\le Y}d\le Y^2\) for \(Y\ge1\),
and the decreasing semicircle sum is bounded by its integral.
The union bound and Lemma~\ref{lem:wheel} give the tail bound.
When its threshold is below one there is no nonzero integral vector;
when it is at least \(M\), the stated probability bound is trivial.
Finally
\[
 \E\log Z_M\le \E(\log Z_M)_+
 \le\int_0^\infty\min(1,Ke^{-2t})\,dt
       =(\log K+1)/2,\quad K=J(q)^2+3.
\]
\end{proof}

For two primes a slightly sharper tail is available:
\[
 \Prb(Z_M>z)\le\min(1,6z^{-2}).                       \tag{2.2}\label{eq:two_tail}
\]
Indeed, if neither denominator of \(u,v\) is a full unit, the two bad-prime
sets must be complementary. Then both \(u+v\) and \(u-v\) have unit
denominators. The smaller Euclidean norm is at most
\(\sqrt{\|u\|^2+\|v\|^2-2|\langle u,v\rangle|}\le\sqrt2 M/\lambda_s\)
by the reduced-basis bounds. If one basis vector is a unit the same
inequality follows directly. Combining with the disk count proves
\eqref{eq:two_tail}.

\begin{proof}[A first uniform bound]
Prime-power bases satisfy the stronger second assertion, proved below.
Suppose \(q\) has at least two primes and \(q\ge10\). If \(q\) is odd,
the units \(1,(q-1)/2,(q+1)/2,q-1\) show \(J(q)\le q/2\).
If \(q=2d\) is even and \(d\ge5\) is odd, the units
\(1,d-2,d+2,2d-1\) give the same conclusion. Thus
\[
 J(q)^2+3\le 7q^2/25,\qquad e(J(q)^2+3)<q^2,
\]
using \(e<3\). Equation~\eqref{eq:wheelmean} gives
\(\mu(B^N)<\log q\le\log B\).
For \(q=6\), \(J(6)=4\); if \(B\ge12\), \(19e<57<144\le B^2\).
For \(B=6\), \eqref{eq:two_tail} gives
\(\mu(M)\le(\log6+1)/2<\log6\).
All constants are independent of \(N\).
\end{proof}


\begin{proposition}[An explicit common correction]\label{prop:uniform-correction}
For all integer bases \(B\ge2\),
\[
 c(B)\le\max(0,w,39/40)<49/50,\qquad
 w=1+\frac{\log(209/300)}{24\log2}.
\]
In particular \(C^*<49/50\) and Theorem~\ref{thm:main}(1) holds.
\end{proposition}
\begin{proof}
The prime-power and two-prime arguments of Sections~\ref{sec:prime}
and~\ref{sec:two} do not use this proposition; they give \(c(B)\le0\)
when \(k=\omega(q)\le2\), where \(q=\rad(B)\).
For the other bases, \eqref{eq:wheelmean} gives
\[
 c(B)\le\frac{\log(e(J(q)^2+3))}{2\log q}.
\]
If \(3\le k\le5\), then \(q\ge30\) and
\(J(q)\le\min(q/2,2048)\).
For \(q\le4096\), using \(e<11/4\),
\[
 \frac{e(J(q)^2+3)}{q^2}
 <\frac{11}{4}\left(\frac14+\frac3{30^2}\right)=\frac{209}{300}.
\]
Put \(\delta=-\log(209/300)>0\). Since
\(2\log q\le24\log2\), the common upper bound is \(c(B)\le w\).
For \(q\ge4096\), the fixed numerator \(\log(e(2048^2+3))\)
and the increasing denominator give the same bound, by evaluating
at \(q=4096\). Moreover
\[
 \delta>\frac{91}{300}+\frac12\left(\frac{91}{300}\right)^2
 =\frac{62881}{180000}>\frac{42}{125},\qquad
 24\log2<\frac{84}{5}.
\]
Thus \(\delta/(24\log2)>1/50\), and \(w<49/50\).

For \(k\ge6\), the product of the first six distinct primes is at least
\(30030\), and every further prime is at least \(17>16\).
Hence \(q\ge30030\,16^{k-6}\), while \(J(q)^2+3\le5\,16^k\).
The elementary enclosures
\(\log30030>10,\ 2/3<\log2<7/10,\ \log5<17/10\) give
\[
 c(B)<\frac{(14/5)k+27/10}{20+(16/3)(k-6)}
       \le\frac{39}{40}<\frac{49}{50}.
\]
The rational function equals \(39/40\) at \(k=6\) and decreases;
its derivative numerator is \(-48\).
For the only additional constant, \(e<11/4\) and
\(11^{10}<30030\,4^{10}\) imply \(\log30030>10\).
Taking the supremum over bases uses the two \emph{common} caps
\(w\) and \(39/40\), rather than merely a pointwise strict inequality.
No optimum is identified.
\end{proof}

\section{Exact central counts and certified logarithms}\label{sec:cert}
Let \(U_m(T)=\#\{1\le d\le T:\gcd(d,m)=1\}\).
For \(T^2\le M\) and a proper divisor \(g\mid M\), put \(m=M/g\) and
\[
 P_g(T)=\sum_{\substack{e\le T/g\\\gcd(e,M)=1}}
     \mu_{\rm Mob}(e)\,
     U_m\!\left(\floor{\frac{T}{ge}}\right)
     U_q\!\left(\floor{\frac Te}\right),\quad q=\rad(M). \tag{3.1}\label{eq:Pg}
\]
Here \(\mu_{\rm Mob}\) denotes the Möbius function.
These count primitive positive witnesses \(a=ga'\), with
\(\gcd(a',m)=1\) and full-unit denominator.
Within either sign no two such primitive fractions coincide modulo \(M\):
their determinant has absolute value at most \(T^2-1<M\).
Opposite signs can collide only when \(af+ed=M\); the possible value \(2M\)
at \(T^2=M\) would force all four components equal to \(T\), which cannot
give a unit denominator. Write \(J_g(T)\) for these collisions in the
gcd-\(g\) state layer. Then
\[
 C_g(T)=2P_g(T)-J_g(T),\qquad
 C(T)=1+\sum_{g\mid M,\ g<M}C_g(T).                   \tag{3.2}\label{eq:central}
\]
The one is the zero state and must not be absorbed into a nonzero layer.

For completeness the collisions are finite exact computations:
enumerate ordered full-unit \(d,f\le T\) with \(\gcd(d,f)=1\),
set \(a_0=Mf^{-1}\bmod d\) with \(0\le a_0<d\), and
\(e_0=(M-fa_0)/d\).
All solutions are
\(a=a_0+kd,\ e=e_0-kf\).
Intersect the two closed integer ranges \(1\le a,e\le T\);
assign surviving pairs to \(\gcd(a,M)\), enforcing primitivity.
This handles endpoints, signs and empty phases exactly.

\subsection{Lower CDFs outside the central range}
For \(M=p^a r^b\), \(p<r\), a nonzero state in layer \(g\) is \(s=gt\),
where \(t\) is a unit modulo \(m=M/g\).
For \(m>2\), choose the centered inverse magnitude \(h\) and its sign
\(\varepsilon\in\{1,-1\}\), so that
\(1\le h\le H=\floor{(m-1)/2}\) and \(th\equiv\varepsilon\pmod m\).
The witness \((\varepsilon g,h)\) is admissible if \(h\) is a full unit.
If \(m\) is a pure power of one prime \(u\) and the missing prime \(v\)
divides \(h\), replacing \(h\) by \(m-h\) repairs the denominator and
changes the numerator sign. Put
\[
 V(x)=\floor{x/v}-\floor{x/(pr)}.
\]
For \(T\ge g,m>2\) the resulting lower CDF in that layer is
\[
 W_g(T)=2U_q(\min(T,H))
 +2\unit_{\{m\ {\rm pure}\}}
  \max(0,V(H)-V(\max(0,m-T-1))).                    \tag{3.3}
\]
For \(T<g\) it is zero. When \(m=2\) the layer contains one state,
with witness height \(g\). There is also the direct-denominator-one bound
\[
 B_g(T)=2U_m(\min(\floor{T/g},H))                 \tag{3.4}
\]
for \(T\ge g,m>2\), with the same exceptional conventions.
Each layer has the deterministic height cap
\[
 K_g=\min\bigl(M/p,\;g\floor{m/2},\;\max(g,D_m)\bigr), \tag{3.5}
\]
where \(D_m=1\) for \(m=2\), \(D_m=H\) when \(m\) has both primes, and
\(D_m=\max(H,m-v)\) when \(m\) is pure.
The first cap follows from the circle lemma below; the others follow
from the displayed witnesses. Bounds from the same layer are combined
by a maximum, not a sum.

\begin{lemma}[Circle cap]\label{lem:circle}
For \(M\) with smallest prime divisor \(p\),
\(R_M(s)\le\max(M/p,p-1)\).
If \(M/p\ge p-1\), the bound is \(M/p\), and equality occurs for
\(s=M/p\). This condition holds for every two-prime modulus and
every prime power of exponent at least two.
\end{lemma}
\begin{proof}
Place \(0,s,\ldots,(p-1)s\) on the circle of circumference \(M\).
If two coincide their index difference supplies numerator zero and a
denominator between one and \(p-1\). Otherwise their smallest cyclic
gap is at most \(M/p\), again with index difference between one and
\(p-1\). That denominator is coprime to \(M\), proving the maximum bound.
For \(s=M/p\), every unit denominator is nonzero modulo \(p\);
its centered product with \(s\) has absolute value at least \(M/p\).
\end{proof}

For \(T^2>2M\), delete nonunit denominators from the Farey sequence of
order \(T\). Adjacent deleted denominators with complementary bad primes
have unit difference; taking that difference repairs one neighbor.
A gap can remain uncovered only next to a bad center \(d<M/T\).
Its two retained neighbors have denominators \(>T-jd\), where \(j=1\)
if \(pr\mid d\), and \(j=2\) otherwise. These are positive because
\(T^2>2M\). Its open length in residue coordinates is at most
\(2(M-Td)/(d(T-jd))\). One extra integer per open hole gives
\[
 S(T)\le
 \sum_{\substack{d\le\floor{(M-1)/T}\\\gcd(d,M)>1}}\varphi(d)
 +\sum_{\substack{d\le\floor{(M-1)/T}\\\gcd(d,M)>1}}
       \frac{2\varphi(d)(M-Td)}{d(T-jd)}.           \tag{3.6}\label{eq:far}
\]
This is an upper survival bound, not an equality.
The strict start \(\floor{\sqrt{2M}}+1\) avoids the removable
equality endpoint. The same Farey repair and length argument underlies
the analytic estimates in the appendix.


\subsection{Pointwise histograms for the initial finite bridges}
\label{sec:early-histograms}
For a two-prime modulus \(M\), the base-six diagonal, sixth-power and
squarefree-prefix batches use a pointwise construction, distinct from
the aggregate CDF calendar. Put
\[
 k=\floor{\sqrt M},\qquad \tau=\floor{\sqrt{2M}}+1,\qquad
 D_1(s)=\max(1,\min(s,M-s)).
\]
The closed central enumeration gives the exact height of every state
with \(R_M(s)\le k\), including zero separately.
Apply the proved far-cover construction at \(\tau\), retaining its
actual full-unit neighbors and phases, and let \(\mathcal H\) be the
union of the open exceptional holes.
Outside \(\mathcal H\), set \(U(s)=\min(\tau,D_1(s))\), since the cover
proves \(R_M(s)\le\tau\).
Inside \(\mathcal H\), start with \(D_1(s)\) and lower \(U(s)\) only
by heights \(\max(|a|,d)\) of explicitly checked pairs
\[
 d>0,\qquad \gcd(d,M)=1,\qquad a\equiv sd\pmod M.
\]
Every such pair is an admissible upper witness; overlaps are handled
by minima. Restore the exact central heights wherever available.
Consequently \(R_M(s)\le U(s)\) pointwise, and its histogram has mass \(M\).

For example, at a reduced exceptional center \(c/v\), put
\(\Delta=Mc-vs\) and \(\varepsilon=\operatorname{sign}(\Delta)\),
with \(\varepsilon=1\) if \(\Delta=0\).
For any positive full-unit \(d\) with \(cd\equiv\varepsilon\pmod v\), the
integer
\[
 a=\varepsilon(M-|\Delta|d)/v
\]
satisfies \(a-sd=M(\varepsilon-cd)/v\).
Thus it is a legal witness; its optimality is neither assumed nor needed.
The implementation checks integrality, congruence and denominator marks.

For \(M=pr\), \(p<r\), cap \(U\) by \(r\), using Lemma~\ref{lem:circle}.
When \(r\le\tau\), that cap suffices off the exact central set.
For the actual histograms in these batches, the separately checked
integer sign premise \(Q<0\), together with \(R_M\le U\) and the outward
logarithm enclosures below, proves the required negative centered mean.
Pointwise domination alone does not imply that sign.
For the diagonal moduli \(M=6^n\), \(2M\) is nonsquare, so the historical
threshold \(\ceil{\sqrt{2M}}\) equals the strict threshold \(\tau\).

\subsection{Integer logarithmic certificates}
A calendar contains all integers up to an initial \(k\), then steps by
\(T\mapsto\ceil{257T/256}\), with every layer cap, the far-tail start,
and the terminal \(M/p\) inserted. At each calendar point combine
\eqref{eq:central} when applicable, the layer witness bounds, and
\eqref{eq:far}; take cumulative maxima. Filling the gaps by the previous
CDF value produces a lower CDF \(\underline C\le C\).
Its increments \(n_h\) give a dominating height histogram:
\[
 \sum_hn_h=M,\qquad \E\log R_M\le M^{-1}\sum_hn_h\log h.
\]
The increments are not claimed to be the true height distribution.

Use outward integer ticks \(L_S(x)\le S\log x\le U_S(x)\) and retain
\[
 Q=2\sum_h n_h U_S(h)-M L_S(M).                    \tag{3.7}\label{eq:Q}
\]
Then \(Q<0\) proves \(\mu(M)<0\), including odd \(M\).
The identity
\[
 \log x=2\sum_{j=0}^{K-1}\frac{z^{2j+1}}{2j+1}
       +\varepsilon,\quad z=(x-1)/(x+1),\quad
 0\le\varepsilon\le
 \frac{2z^{2K+1}}{(2K+1)(1-z^2)}
\]
for rational \(x\ge1\), with range reduction by powers of two, supplies
exact rational enclosures. No floating logarithm or assertion-sensitive
test is used for the sign. The independent checker uses a different
series depth and tick scale.
It independently reconstructs the initial boundary, terminal boundary
and relevant boundary gcd strata. It consumes the intervening central
arrays; the central lemma and producer are therefore essential parts
of the proof. It is not a complete independent recount of every
intermediate CDF value.

\section{Prime powers}\label{sec:prime}
For odd \(M=p^n\) let \(a=\sqrt M\), \(K=\floor{a/\sqrt2}\).
The central count is
\[
 C(T)=1+2P_p(T),\qquad
 P_p(T)=\#\{u,d\le T:\gcd(u,d)=1,\ p\nmid d\},\quad T\le K.
\]
Indeed \(2T^2<M\), so opposite primitive witnesses do not collide.
Before primitivity there are at least \((1-1/p)T^2\) pairs.
The divisor union bound gives
\[
 P_p(T)\ge(1-y)\{2-(1-y^2)\zeta(2)\}T^2
            -(1-y)T\log T,\qquad y=1/p.
\]
To see the error explicitly, a common divisor \(d\ge2\) must satisfy
\(p\nmid d\), and
\(\floor{x}(\floor{x}-\floor{x/p})\le(1-y)(x^2+x)\).
Sum over \(d\), use
\(\sum_{d\ge2,p\nmid d}d^{-2}=(1-y^2)\zeta(2)-1\)
and \(\sum_{d=2}^T d^{-1}\le\log T\).
The elementary enclosure \(\zeta(2)<5/3\) follows by summing five terms
and bounding the rest by an integral. Consequently
\[
 C(T)\ge cT^2-dT\log T,\quad
 (c,d)=
 \begin{cases}(56/81,4/3),&p=3,\\(5/8,2),&p\ge5.\end{cases}              \tag{4.1}
\]

The exact compensation identity is
\[
 \mu(M)=-D+P_+,\quad
 D=M^{-1}\int_1^a C(\floor x)\frac{dx}{x},\quad
 P_+=M^{-1}\int_a^\infty S(\floor x)\frac{dx}{x}.     \tag{4.2}\label{eq:comp}
\]
Retaining the plateau \(C(K)\) up to \(a\) yields
\[
 D\ge c(K^2/a^2)(1/2+\log(a/K))
      -2cK/a^2-d(K\log K/a^2)(1+\log(a/K)).          \tag{4.3}
\]
For \(a\ge140\), \(K/a\ge11/16\), \(K/a<3/4\),
\(K^2/a^2\ge1/2-\sqrt2/a\), and
\(\log2/2\le\log(a/K)\le5/11\).
Using \(2/3\le\log2\le1\) and \(\sqrt2<3/2\) gives
\[
 D\ge
 \begin{cases}
 70/243-[56/27+(16/11)\log a]/a,&p=3,\\
 25/96-[15/8+(24/11)\log a]/a,&p\ge5.
 \end{cases}                                      \tag{4.4}
\]

For integer \(T\ge a\), the holes next to reduced Farey centers with
\(p\mid d,d<M/T\) have retained neighbors \(b_L,b_R>T-d\).
Their open lengths are at most \(2(M-dT)/(d(T-d))\).
Let \(\kappa=(p-1)/p^2\), \(t=T/a\), and
\[
 F(t)=\sum_{j\ge1}\frac{t^{-2j}}{j(j+1)}
     =1+(t^2-1)\log(1-t^{-2}),\quad F(1)=1.
\]
Summation of lengths and one integer per hole, using
\(\varphi(pk)\le(p-1)k\), gives
\[
 S(T)/M\le2\kappa F(t)+\frac{\kappa a^2}{2T^2}
                  +\frac{1-1/p}{2T}.              \tag{4.5}
\]
One can check the rounding term directly: the hole count is bounded by
\((p-1)L(L+1)/2,\ L=\floor{M/(pT)}\).
The decreasing length summand is integrated with spacing \(p\).
Since \(F\) decreases and
\(\int_1^\infty F(t)\,dt/t=(\zeta(2)-1)/2\), integrating the initial
interval \([a,\ceil a]\) and retaining the first summand gives
\[
 P_+\le
 \begin{cases}
 11/54+17/(9a)+1/(3a^2),&p=3,\\
 11/75+19/(10a)+1/(2a^2),&p\ge5.
 \end{cases}
\]
Thus
\[
 \mu(M)\le
 \begin{cases}
 -41/486+[107/27+(16/11)\log a]/a+1/(3a^2),&p=3,\\
 -91/800+[151/40+(24/11)\log a]/a+1/(2a^2),&p\ge5.
 \end{cases}                                      \tag{4.6}\label{eq:primeupper}
\]
At \(a=140\), using \(\log140<5\), both upper bounds are negative:
the positive margins are at least \(214289/52390800\) and
\(3811/431200\), respectively. All error terms decrease thereafter.

For \(M=p\ge5\), choose \(k=\floor{\sqrt p}\). Among the
\((k+1)^2>p\) pairs \(0\le u,v\le k\), two have the same \(u-sv\);
their difference has nonzero denominator at most \(k<p\).
Thus \(R_p\le k<\sqrt p\).
For \(M=p^2\), the circle argument gives \(R_M\le p=\sqrt M\).
Only the 21 entries in Table~\ref{tab:prime} remain below \(19600\).
Here \(H=\sum_sR_M(s)\) and the exact inequality \(H^2<M^3\)
implies \(\mu(M)<0\) by Jensen.
The companion computes these minima without an assumed mean bound.

\begin{table}[ht]
\centering
\begin{tabular}{lll}\toprule
Prime & Exponents \(n\) & Sums \(H\), in order\\\midrule
3 &1--8&3,17,105,577,3107,16527,86977,455939\\
5 &3--6&1009,11729,134011,1513793\\
7 &3--5&4447,84975,1595511\\
11&3--4&32803,1220493\\
13&3&68697\\
17&3&225957\\
19&3&370511\\
23&3&866951\\\bottomrule
\end{tabular}
\caption{The finite odd-prime bridge.}\label{tab:prime}
\end{table}
This proves Theorem~\ref{thm:main}(2) for odd prime powers.
If \(B=p^b\), replace \(n\) by \(bN\); no new finite cases occur.
For \(p=2\), use the exact theorem of \cite{binary}.

\section{The two-prime theorem}\label{sec:two}
The appendix gives the full parametric compensation inequality and its
refinements. Its proof uses the exact central formula, explicit errors
for the unit-counting function, and the Farey survival estimate
\eqref{eq:far}. It is an inequality at each finite \(M\);
no limiting approximation is substituted for a finite bound.
In particular, for \(M=p^a r^b\), it proves negativity above:
\[
 \begin{array}{c|c}
 \hbox{prime factors}&\hbox{sufficient modulus}\\\hline
 2,3 &6^9\\
 2,r,\ r\ge5 &9\,000\,000\\
 3\le p<r &2^{22}
 \end{array}                                      \tag{5.1}
\]
The sharper uniform cutoffs, apart from the four radicals
\(6,10,14,15\), are
\[
 M\ge C(p),\quad
 C(2)=384^2,\quad C(3)=512^2,\quad C(p)=600^2\ (p\ge5).                 \tag{5.2}\label{eq:cutoffs}
\]
Appendix~\ref{ap:finite-two-prime} specifies all finite domains and reductions.
The certificates are organized by domains of moduli, not assumed
experimental monotonicity of means.

\begin{proof}[Proof of Theorem~\ref{thm:main}(3)]
First the analytic bounds and the finite bridges settle the four
exceptional radicals \(6,10,14,15\) at every pair of positive valuations.
The squarefree two-prime moduli are settled by a single-mark reduction,
with 4596 smaller moduli certified explicitly.
For all bases, the analytic cutoffs leave three sixth powers, three
fourth powers, 35 cubes and 99 squares, as detailed in the appendix.
These imply the assertion for every length \(N\ge2\).
The stronger length cutoffs are reindexings of the same modulus, not
new hypotheses on states.

For \(N=1\), write \(M=p^a r^b\). If \(\gcd(a,b)\ge2\), it is already
a higher power of a two-prime base. Remove squarefree moduli, the four
settled radicals, and the moduli above \eqref{eq:cutoffs}.
The resulting exact inventory has 19239 profiles.
A larger-prime-exponent-one lemma pays 7857 profiles.
The five explicitly proved analytic bounds pay another 6972.
The remaining 4410 profiles split into the three disjoint certificate
domains of sizes 1087,1422 and1901. Each retained integer \(Q\) in
\eqref{eq:Q} is strictly negative.
All inventories, prime tests, powers, exclusions and sign comparisons
are finite integer operations. The domain identities and mass totals
are checked in the companion. There is no unexamined valuation class.
\end{proof}

\section{A lower estimate and positive corrections}
For squarefree \(q\), put
\[
 D_q=2^{\omega(q)},\quad
 \rho_q=\prod_{p\mid q}\frac p{p+1},\quad
 \alpha_q=\frac{2\rho_q}{\zeta(2)},\quad
 \beta_q=-\frac{1+\log\alpha_q}{2}.
\]

\begin{lemma}\label{lem:lower}
With \(q=\rad(M)\) and \(L=\sqrt{M/\alpha_q}\),
\[
 \mu(M)\ge
 \beta_q-
 \left(\frac{2(D_q+3)}{\sqrt{\alpha_q M}}+\frac1M\right)\log L.
                                                               \tag{6.1}\label{eq:lower}
\]
\end{lemma}
\begin{proof}
Let \(P_q(T)\) count positive primitive pairs \(a,d\le T\) with
\(\gcd(d,q)=1\). Möbius inversion gives
\[
 P_q(T)=\sum_{\substack{e\le T\\\gcd(e,q)=1}}
       \mu_{\rm Mob}(e)\floor{T/e}U_q(\floor{T/e}).
\]
Inclusion--exclusion gives \(|U_q(x)-(\varphi(q)/q)x|\le D_q\).
Expanding each product, summing \(e^{-1}\le1+\log T\), and bounding
the tail of \(\sum e^{-2}\) by \(1/T\), yields
\[
 \left|P_q(T)-\frac{\rho_q}{\zeta(2)}T^2\right|
       \le(D_q+3)T(1+\log T).                       \tag{6.2}
\]
The main coefficient follows from
\((\varphi(q)/q)\sum_{\gcd(e,q)=1}\mu_{\rm Mob}(e)e^{-2}
=\rho_q/\zeta(2)\). These bounds also cover \(T=1\).
Every admissible fraction can be reduced to a primitive one, so
\(C(T)\le1+2P_q(T)\) by the union bound; injectivity is not needed.
Use the resulting upper bound for \(C(\floor x)\) on \(1\le x\le L\)
in \(\E\log R_M=\int_1^\infty(1-C(\floor x)/M)\,dx/x\), and discard
the nonnegative remainder. Integrating gives \(\beta_q\) as the main
centered term. The error integral is at most
\(2(D_q+3)L\log L/M+\log L/M\).
The omitted lower endpoint main term is nonnegative, giving
\eqref{eq:lower}.
\end{proof}
For fixed \(B\), this bound and \eqref{eq:wheelmean} yield
\(E_{B,N}/N\to1/2\), consistently with the antecedent in \cite{GK}.

\subsection{The exact finite obstruction}
The direct height histogram for \(M=900\) has total mass 900 and height
sum 34007. If
\(P=\prod_{s=0}^{899}R_{900}(s)\), integer multiplication gives
\[
             2\cdot30^{900}\le P<3\cdot30^{900}.     \tag{6.3}\label{eq:counter}
\]
The second bound is recorded for reproducibility, not needed for the sign.
Hence \(\E\log_{30}R_{900}=\log_{30}P/900>1\).
An independent algorithm enumerates admissible denominators and centered
numerators directly. Replacing full-unit denominators by denominators
coprime to only one prime changes many minima; that wrong model is not
used. The decimal \(1.0002564533304775\ldots\) is descriptive;
the proof is the integer comparison \eqref{eq:counter}.
It refutes a possible strengthening, not a conjecture stated in
Question 3.2.

\subsection{An eventually positive family}
For \(q_0=2\cdot3\cdot5\cdot7\cdot11\cdot13\cdot17=510510\),
\[
 \rho_{q_0}=12155/41472.
\]
For \(M\ge q_0^2\), the main lower term in \eqref{eq:lower} satisfies
\(\beta_{q_0}>7/500\), and its error is \(<13/1000\).
We give rational checks to make this explicit.
The elementary enclosures \(e<87/32\) and \(\zeta(2)>41/25\)
imply
\[
 e\alpha_{q_0}<1-\frac{256169}{9068544},\qquad
 \beta_{q_0}>\frac{256169}{18137088}>\frac7{500}.
\]
The first follows from a factorial-series upper bound for \(e\),
the second from a partial sum for \(\zeta(2)\).
Also \(\zeta(2)<5/3\) gives
\(\alpha_{q_0}>2431/6912>7/20>(59/100)^2\).
At \(M=q_0^2\), \(\log L<14\) and \(q_0>510000\), so
\[
 \left(\frac{2(128+3)}{\sqrt{\alpha_{q_0} M}}+\frac1M\right)\log L
 <\left(\frac{262}{(59/100)510000}
               +\frac1{510000^2}\right)14<\frac{13}{1000}.
\]
The functions \(\log(\sqrt{M/\alpha})/\sqrt M\) and
\(\log(\sqrt{M/\alpha})/M\) decrease on this range.
Thus \(\mu(M)>1/1000\), proving the family in
Theorem~\ref{thm:positive}. No minimal radical or minimal cutoff is claimed.
For radicals containing \(q_0\), an eventually positive conclusion also
follows from \eqref{eq:lower} with that radical's own constants.

\section{The global additive optimum}
\begin{proof}[Proof of Theorem~\ref{thm:optimal}]
Let \(k=\omega(q)\). Since the \(i\)-th smallest distinct prime is at
least \(i+1\), \(q\ge(k+1)!\).
Lemma~\ref{lem:wheel} therefore gives \(J(q)=q^{o(1)}\) uniformly
over squarefree radicals: for \(k\to\infty\),
\[
 \frac{\log J(q)}{\log q}
 \le\frac{\log2+k\log4}{\log((k+1)!)}\longrightarrow0;
\]
for bounded \(k\) and \(q\to\infty\), the numerator is bounded.
Equation~\eqref{eq:wheelmean} gives
\[
 c_+(B)\le\frac{\log(J(\rad B)^2+3)+1}{2\log B}.
\]
If the radical tends to infinity this goes to zero, since
\(\log q\le\log B\). If radicals stay bounded their numerator is bounded
and \(\log B\to\infty\). This proves \(c_+(B)\to0\).

The witness \(c(30)\ge E_{30,2}-1>0\) now excludes all sufficiently large
bases from the global supremum. Hence only finitely many bases remain.
Proposition~\ref{prop:uniform-correction} supplies the common cap
\(C^*\le\max(0,w,39/40)<49/50\). Together with the positive witness,
this gives \(0<C^*<49/50\).
A finite maximum of the quantities \(c(B)\) is still a maximum of
suprema over infinitely many lengths; finite-length attainment is not
part of the conclusion.
\end{proof}
This elementary argument requires no hypothesis about extremal
Jacobsthal gaps or prime distributions. It is not a practical algorithm
for locating the maximizing base.


\begin{corollary}\label{cor:extremal-base}
If \(c(B)=C^*\), then \(B\) has at least three distinct prime factors and
is not a perfect power. Equivalently, the gcd of its positive prime
valuations is one.
\end{corollary}
\begin{proof}
The classes with one or two primes have \(c(B)\le0<C^*\).
For \(B=A^e\), \(e\ge2\), equality of modulus and witness spaces gives
\[
 E_{A^e,N}-N/2=\frac1e(E_{A,eN}-eN/2),\qquad
 c(A^e)\le c(A)/e\le C^*/e<C^*.
\]
The supremum here is over a subset of the lengths for \(A\); no
maximizing finite length is asserted. Non-perfect-power does not
mean squarefree.
\end{proof}

\begin{corollary}[A classical Jacobsthal sharpening]\label{cor:kanold}
Using Kanold's bound \(J(q)\le2^{\omega(q)}\), as reported by
Costello--Watts~\cite[Introduction, pp.~1--2]{CostelloWatts}, one has
\[
 0<C^*\le\frac{\log67+1}{2\log30}
       <\frac{5037}{6400}<\frac45.
\]
In particular, \(E_{B,N}<N/2+4/5\) for all \(B\ge2,N\ge1\).
\end{corollary}
\begin{proof}
The maximum cyclic gap \(J(q)\) equals the least length of a consecutive
integer interval guaranteed to contain a \(q\)-unit: a gap of length \(G\)
has \(G-1\) intervening nonunits. Thus the cited bound uses precisely our
convention, and applies to every radical, not only primorials.

Put \(q=\rad(B)\) and \(k=\omega(q)\). For \(k\le2\),
Theorem~\ref{thm:main}(2)--(3) gives \(c(B)\le0\).
For \(k\ge3\), equation~\eqref{eq:wheelmean}, uniformly in \(N\), gives
\[
 c(B)\le\frac{\log(4^k+3)+1}{2\log q}.
\]
The first three distinct prime divisors have product at least \(30\),
and each further prime is at least \(7\), hence at least \(4\).
Consequently \(q\ge30\,4^{k-3}\). Also
\(\log(4^k+3)\le k\log4+\log(67/64)\), so
\[
 c(B)\le F(k):=
 \frac{k\log4+1+\log(67/64)}{2\log30+2(k-3)\log4}.
\]
The denominator is positive for real \(k\ge3\); the numerator of its
derivative is \(2\log4(\log30-\log67-1)<0\).
Thus \(F(k)\le F(3)=(\log67+1)/(2\log30)\).
This is a common cap before either supremum over \(N\) or \(B\).

For a rational enclosure, Appendix~\ref{app:constants} gives
\(\log2<7/10\) and \(e<11/4\). Together with \(\log(1+x)<x\),
\[
 \log67+1=6\log2+\log(67/64)+1
 <\frac{21}{5}+\frac3{64}+1=\frac{1679}{320}.
\]
Since \(11^{10}=25937424601<28311552000=27000\,4^{10}\),
one has \(e^{10}<30^3\), hence \(\log30>10/3\). Therefore
\[
 \frac{\log67+1}{2\log30}
 <\frac{1679/320}{20/3}=\frac{5037}{6400}
 <\frac{5120}{6400}=\frac45.
\]
Positivity follows from Theorem~\ref{thm:positive}.
\end{proof}
This sharpening imports the cited classical gap bound; the self-contained
\(49/50\) argument above remains valid independently.
The argument does not determine the optimal correction or its extremizers.

\section{Reproducibility and limitations}\label{sec:repro}
The repository companion contains the original finite mathematical premises,
portable reconstruction programs, exact logarithmic sums, complete domain
lists and a provenance map. The original default checks remain a short
retained-evidence check, a fixed defining oracle and a four-case squarefree
reconstruction. They are distinct from the complete reconstruction route
provided by \texttt{scripts/reconstruct\_all.py}. The reproduction guide
specifies the commands, domains and limits of both routes.

The complete source-replay chain passed all 37 stages in each of normal
and optimized Python modes. The squarefree prefix, which originally
retained only digests and signs, now has all 4596 reconstructed upper
histograms and their consumed logarithmic ledgers. These histograms cover
41585055 original envelope entries. A separate K32 logarithm implementation
checked all 705257 consumed tick entries, involving 6929 distinct integers,
and reconstructed every negative outward sum and its exact normalization.
The retained full artifacts are indexed by
\texttt{companion/reconstruction/index.json};
\texttt{scripts/check\_complete\_evidence.py} checks their declared complete
coverage, byte identities and the squarefree logarithmic certificates.

An additional distinct central checker reconstructs all integer thresholds
\(1\le T\le k=\floor{\sqrt M}\) and every proper original gcd stratum
for all nine aggregate batches, including their consumed validation rows.
For each full-unit denominator \(1\le d\le k\), it enumerates every
\(1\le a\le k\) and both signed residues \(\pm ad^{-1}\pmod M\),
retaining the smallest \(\max(a,d)\) at each state. There is no
\(\gcd(a,d)=1\) filter: nonprimitive pairs are legal witnesses too.
Zero is one separate atom; coincident signed witnesses update the same
state. The resulting central height bins give every stratum prefix.
All source-produced central arrays are compared at every threshold,
including explicit zero counts for omitted strata above \(k\).
The original endpoint-only checker descriptions elsewhere in the paper
refer to the earlier retained-data checks, rather than this additional
complete central comparison.

The same direct central enumeration also checks every central histogram
prefix for all 4596 squarefree cases, the eight base-six cases and the
three sixth-power cases. All nine aggregate and all three pointwise-domain
checks completed in both normal and optimized modes. The pointwise
histograms provide global prefix comparison targets; they do not retain
separate source stratum arrays. These checks compute exact central minima,
not complete true minimum vectors in the far range. Far-envelope domination,
the infinite estimates and the reduction arguments still use the proved
mathematics and the source producer. The attributed binary mean theorem
remains an import; this reconstruction does not repeat its earlier paper.

The paper proves mean bounds for its full-unit rational-height model.
It does not determine the optimal correction, maximizing lengths or a
sign classification for all bases with at least three primes.
It makes no assertion about the Collatz conjecture. The classical tools
and earlier binary theorem are attributed, and no universal priority
claim is made.

\paragraph{Data and code availability.}
All author-generated integer certificates and programs needed for the
finite proof dependencies accompany the paper in the repository.
The manuscript is licensed under CC BY 4.0; original code, certificates
and repository documentation are licensed under MIT. Cited publications
retain their own rights and are not bundled.

\paragraph{AI-use disclosure.}
OpenAI GPT models assisted with mathematical exploration, proof criticism,
literature searches, manuscript editing, implementation and reproducibility
checks. Several AI review passes shared source exposure. These do not
constitute independent human peer review, formal proof verification or an
originality certificate. The author retains responsibility for the arguments,
source attributions and accompanying artifacts.

\appendix
\section{A finite compensation inequality for two-prime moduli}
\label{ap:two-prime-compensation}

Let \(M=p^a r^b\), where \(p<r\) are distinct primes and \(a,b\geq1\).
Throughout this appendix the average is over all \(M\) residues, including zero,
and
\[
 R_M(s)=\min_{\substack{d>0,\ (d,M)=1\\ A\equiv sd\pmod M}}
          \max(|A|,d),\qquad
 \mu(M)=\frac1M\sum_{s\bmod M}\log R_M(s)-\frac12\log M .
\]
All logarithms are natural. Write
\[
 \begin{gathered}
 A_0=\sqrt M,\quad D=pr,\quad u=(1-p^{-1})(1-r^{-1}),\\
 \rho=\frac{pr}{(p+1)(r+1)},\quad
 \alpha=\zeta(2)^{-1},\quad \ell=\log2.
 \end{gathered}
\]
\[
 \kappa=\alpha\rho,\qquad
 \theta=\alpha(1-\rho),\qquad \beta=3-2\sqrt2 .
\]
The symbol \(A_0\) denotes a real scale, not a numerator.

\subsection{Unit discrepancy and primitive counting}
\label{ap:unit-discrepancy}
Define
\[
 U(x)=\lfloor x\rfloor-\lfloor x/p\rfloor-\lfloor x/r\rfloor
                         +\lfloor x/D\rfloor,\qquad e(n)=U(n)-un.
\]
Let \(h=\max_{0\leq n<D}e(n)\), \(m=\min_{0\leq n<D}e(n)\), and put
\[
 K=h-m+u,\qquad W=\max(h,-m+2u),\qquad H_0=u+h.
\]
Then \(h\geq0\), \(m\leq0\), and \(H_0\geq1\), the last inequality following
from \(e(1)=1-u\). The real discrepancy lies in \([m-u,h]\).
Every closed interval of length \(v\) in the positive half-line therefore
contains at most \(uv+K\) positive \(D\)-units: for its integer endpoints
\(i=\lceil x\rceil\), \(j=\lfloor y\rfloor\), use
\(U(j)-U(i-1)\) and \(j-i\leq y-x\).
This includes both endpoint atoms and zero-width intervals.
Partial summation gives
\[
 \sum_{\substack{d\leq Q\\(d,D)=1}}\frac1d\leq H_0+u\log Q
 \quad(Q\geq1).
\]
For \(x=n+t\), \(0\leq t<1\),
\[
 nU(n)-ux^2=ne(n)-u(2nt+t^2),\qquad
 |nU(n)-ux^2|\leq Wx.
\]
For the lower bound one uses
\[
 (-m+2u)x-\{-nm+u(2nt+t^2)\}
 =2un(1-t)+(-m+2u)t-ut^2\geq0;
\]
the upper bound follows from \(h\geq0\).

Let \(P(T)\) count pairs \(1\leq A,d\leq T\) with
\((A,d)=1\) and \((d,D)=1\). For integer \(T\geq1\), cancelling content
only at \(D\)-units yields
\[
 P(T)=\sum_{\substack{g\leq T\\(g,D)=1}}
       \mu_{\rm Mob}(g)\lfloor T/g\rfloor U(\lfloor T/g\rfloor).
\]
The Euler product
\[
 u\sum_{(g,D)=1}\frac{\mu_{\rm Mob}(g)}{g^2}
   =\frac{u\alpha}{(1-p^{-2})(1-r^{-2})}=\kappa
\]
and the integer-cutoff bound
\(\sum_{g>T}g^{-2}\leq T^{-1}\) prove
\begin{equation}
 |P(T)-\kappa T^2|\leq T(b_0+b_1\log T),
 \label{ap:primitive-error}
\end{equation}
with either of the \emph{coherent} choices
\[
 (b_0,b_1)=(W+u,W)
 \quad\text{or}\quad
 (b_0,b_1)=(WH_0+u,Wu).
\]
Both have \(b_0\geq b_1\). The elementary bounds \(h<2\), \(m>-1\)
also permit the coarse pair \((b_0,b_1)=(4,3)\) and \(K=4\):
here \(|nU(n)-ux^2|\leq3x\), and the unmarked harmonic sum is
at most \(1+\log T\).

For later sharpening we need the stronger uniform integer discrepancy
\begin{equation}
 -u\leq e(n)\leq2u.
 \label{ap:refined-discrepancy}
\end{equation}
It suffices to work in one period. At units \(e(n)-e(n-1)=1-u>0\),
so a minimum occurs at zero or a nonunit. If \(n=pt<D\), write
\(pt=jr+s\), with \(1\leq t<r\), \(0\leq j\leq p-1\), \(1\leq s<r\).
Then \(e(pt)=(s-t)/r\), and
\[
 t-s=jr-(p-1)t
 \leq\frac{jr-(p-1)}p
 \leq\frac{(p-1)(r-1)}p=ur.
\]
Interchanging the primes treats multiples of \(r\).
Reflection gives
\(e(D-n)=\mathbf1_{(n,D)=1}-e(n)\) for \(0<n<D\).
At nonunits this implies \(e(n)\leq u\).
At a unit \(t\), the lower bound for \(e(t-1)\) implies
\(e(t)\geq1-2u\); reflection then gives \(e(n)\leq2u\).
Thus \(K\leq4u\), \(W\leq3u\). For the second coherent parameter choice,
\begin{equation}
 b_0\leq9u^2+u,\quad b_1\leq3u^2,\quad
 b_0-b_1=Wh+u\leq6u^2+u .
 \label{ap:refined-parameters}
\end{equation}

\subsection{Central collisions with closed endpoints}
\label{ap:central-collisions}
For integer \(T^2\leq M\), let \(C(T)=\#\{s:R_M(s)\leq T\}\).
Every nonzero short witness can be cancelled to a primitive one because its
content is a unit modulo \(M\).
Two primitive witnesses with the same numerator sign represent different
states unless they are identical: a nonzero determinant would be a multiple
of \(M\), while its magnitude is \(<M\). At \(T^2=M\), equality in this
bound would require a denominator \(T\), which is not a unit.
Opposite signs can meet only when
\[
 Ad+Eb=M,\qquad 1\leq A,E,b,d\leq T,\quad (bd,D)=1.
\]
The value \(2M\) is impossible for the same nonunit-endpoint reason.
There is at most one witness of either sign at a given state.
Neither a nonzero short witness nor its negative can be the zero state.
Consequently
\begin{equation}
 C(T)=1+2P(T)-J(T),
 \label{ap:central-cdf}
\end{equation}
where \(J(T)\) counts the ordered positive/negative collision pairs above;
no further factor two is inserted.

A common divisor of \(b,d\) divides \(M\) and is coprime to \(M\), so
\((b,d)=1\). For fixed \(b,d\), the interval of progression parameters has
width at most
\[
 \frac{T(b+d)-M}{bd}<1
\]
when nonempty. Indeed
\(T(b+d)-M-bd\leq-(T-b)(T-d)\), with equality requiring \(M=T^2\)
and one nonunit denominator \(T\).
Hence at most one collision occurs per ordered denominator pair.
Put \(L_T=(2T-M/T)_+\). The deficits \(T-b,T-d\) lie in the closed triangle
\(x+y\leq L_T\).
If \(V(v)\) counts admissible deficits up to \(v\), then \(V(v)\leq uv+K\).
Summing the linear weights as a Stieltjes integral gives
\begin{equation}
 J(T)\leq \frac{u^2L_T^2}{2}+2uK L_T+K^2
 \quad(2T^2>M).
 \label{ap:collision-upper}
\end{equation}
There are no collisions for \(2T^2<M\). If \(2T^2=M\), equality forces
\(b=d=T\); the odd prime divisor of \(M\) divides \(T\), so these are
not units. Thus \(J(T)=0\) also at this boundary.

\subsection{The global far cover}
\label{ap:far-cover}
For integer \(T^2\geq2M\), consider the complete Farey sequence of order \(T\).
A unit-denominator center \(c/d\) covers the closed interval
\([c/d-T/(Md),c/d+T/(Md)]\).
Consecutive centers \(a/q<c/d\) have \(cq-ad=1\), \((q,d)=1\), and \(q+d>T\).
If both denominators are units their intervals overlap.
If only one is a unit, an uncovered point forces the other denominator to
be \(<M/T\).

If both are nonunits, their two prime obstructions are complementary.
Hence \(q-d\) is a unit modulo \(D\).
On the Farey arc the affine residual
\[
 (q-d)x-(a-c)
\]
has values \(1/q\) and \(1/d\) at its endpoints.
If \(q,d\geq M/T\), its absolute value multiplied by \(M\) is at most \(T\),
and the vector difference is a genuine witness of denominator
\(0<|q-d|\leq T\) on the whole arc. A witness center outside that arc
is allowed; no isolation of deleted centers is presumed.

It follows that every uncovered point belongs to an arc incident to a
deleted reduced center \(a/q\), \(0<a<q\), with \(q<M/T\leq T/2\).
A complete Farey neighbor \(c/d\) has \(d>T-q>q\).
If \(d\) is also a nonunit, replace it by
\((c-a)/(d-q)\). Its denominator is a unit, lies in \((T-2q,T]\),
and has determinant one with \(a/q\) on the same side.
If \(D\mid q\), the original neighbors are already units and lie in
\((T-q,T]\).
The uncovered residues are therefore contained in the union of open holes
\[
 \left(\frac{Ma}{q}-\frac{M-qT}{qb_L},
       \frac{Ma}{q}+\frac{M-qT}{qb_R}\right),\qquad
 b_L,b_R>T-2q .
\]
An open interval of length \(v\) contains at most \(v+1\) integers.
Overlapping holes may be counted repeatedly in an upper bound. Consequently
\begin{equation}
 M-C(T)\leq
 \sum_{\substack{q<M/T\\(q,D)>1}}\varphi(q)
 \left\{\frac{2(M-qT)}{q(T-2q)}+1\right\}.
 \label{ap:far-holes}
\end{equation}
The \(+1\) allowances, positive denominator marks and open endpoints are
essential. The displayed cover holds also at \(T^2=2M\), but a later
partial summation must start strictly above that threshold.

\subsection{Deleted-center density and its real-cutoff error}
\label{ap:deleted-density}
Put \(A(Q)=\sum_{q\leq Q,\,(q,D)>1}\varphi(q)/q\), and define
\[
 v=2-p^{-1}-r^{-1},\qquad
 c_0=\frac{p-1}{p^2}+\frac{r-1}{r^2}-\frac{D-1}{D^2}.
\]
Expanding \(\varphi(q)/q=\sum_{d\mid q}\mu_{\rm Mob}(d)/d\) and separating its
\(p,r\) parts yields the exact identity
\[
 A(Q)=\sum_{\substack{d\leq Q\\(d,D)=1}}\frac{\mu_{\rm Mob}(d)}d
 \left\{(1-p^{-1})\left\lfloor\frac Q{pd}\right\rfloor
 +(1-r^{-1})\left\lfloor\frac Q{rd}\right\rfloor
 -(1-D^{-1})\left\lfloor\frac Q{Dd}\right\rfloor\right\}.
\]
The linear coefficient in braces is \(c_0Q/d\), and their error is
at most \(v\) in absolute value.
Since \(c_0=(1-p^{-2})(1-r^{-2})-u\), the leading density is \(\theta\).
For real \(Q\geq1\), not merely integer \(Q\),
\[
 Q\sum_{d>Q}d^{-2}\leq Q/\lfloor Q\rfloor<2 .
\]
Thus
\begin{equation}
 |A(Q)-\theta Q|\leq E(Q)=E_0+E_1\log Q
 \label{ap:deleted-error}
\end{equation}
with either \(E_0=v+2c_0,E_1=v\), or \(E_0=vH_0+2c_0,E_1=vu\).
The coarse admissible choice is \(E(Q)=3+2\log Q\), since
\(c_0\leq17/36\). For partial summation on \(0\leq Q<1\), take \(E(Q)=E_0\):
then \(A(Q)=0\), \(\theta Q\leq\theta<1/3\), and \(E_0\geq vu\geq7/18\).
A second partial summation gives
\begin{equation}
 B(Q):=\sum_{\substack{q\leq Q\\(q,D)>1}}\varphi(q)
 \leq\frac{\theta Q^2}{2}+2QE(Q).
 \label{ap:integer-hole-allowance}
\end{equation}

For \(T^2>2M\), write \(Q=M/T\) and
\(f(x)=(M-Tx)/(T-2x)\). It decreases from \(M/T\) to zero at \(Q\).
Equations \eqref{ap:far-holes}--\eqref{ap:integer-hole-allowance} imply
\[
 \frac{M-C(T)}M\leq V(T/A_0)+\frac{4E(Q)}T,
 \qquad
 V(t)=\theta F(t/\sqrt2)+\frac{\theta}{2t^2},
\]
where
\[
 F(z)=\sum_{n\geq1}\frac{z^{-2n}}{n(n+1)},\qquad F(1)=1.
\]
Indeed \(2\theta M^{-1}\int_0^Q f(x)\,dx=\theta F(T/(\sqrt2A_0))\);
the weighted discrepancy costs \(2E(Q)/T\), and the integer-hole
allowance costs another \(2E(Q)/T\).
The function \(V\) is decreasing, \(V(\sqrt2)=5\theta/4\), and termwise
integration of its positive series gives
\begin{equation}
 \int_{\sqrt2}^{\infty}V(t)\frac{dt}{t}
 =\theta\left\{\frac{\zeta(2)-1}{2}+\frac18\right\}
 =\frac{1-\rho}{2}-\frac{3\alpha(1-\rho)}8 .
 \label{ap:far-integral}
\end{equation}

\subsection{The complete finite compensation}
\label{ap:finite-balance}
Denominator one bounds every height by \(\max(1,\lfloor M/2\rfloor)\).
The layer-cake identity is
\[
 \mu(M)=-D_c+P_{\rm mid}+P_{\rm far},
\]
with
\[
 D_c=\frac1M\int_1^{A_0}C(\lfloor x\rfloor)\frac{dx}{x},\qquad
 P_{\rm mid}=\frac1M\int_{A_0}^{\sqrt2A_0}
                  [M-C(\lfloor x\rfloor)]\frac{dx}{x}.
\]
The floor lemma
\[
 \int_1^{A_0}\lfloor x\rfloor^2\frac{dx}{x}\geq A_0^2/2-A_0
\]
follows by summing
\(\log(1+1/n)\geq1/n-1/(2n^2)\) and using
\(\log((k+t)/k)\geq t/(k+t)\) on the last interval.
The residual comparison is at least \(1/2+t-3t^2/2\geq0\) for \(0\leq t<1\).
Integration of \eqref{ap:primitive-error} and
\eqref{ap:collision-upper}, retaining the closed collision boundary, gives
\[
 D_c\geq\kappa-u^2(3/4-\ell)
 -\frac{2\kappa+2(b_0-b_1)+2b_1\log A_0+2uK\beta}{A_0}
 -\frac{K^2\ell}{2A_0^2}.
\]
Here the two collision integrals are
\[
 \frac1M\int_{A_0/\sqrt2}^{A_0}
       (2x-M/x)^2\frac{dx}{x}=3/2-2\ell,\qquad
 \int_{A_0/\sqrt2}^{A_0}(2x-M/x)\frac{dx}{x}=A_0\beta.
\]
The favorable central zero and favorable primitive-error endpoint
\(2(b_0-b_1)/A_0^2\) have only been discarded.
At \(k=\lfloor A_0\rfloor\), \(k^2\geq A_0^2-2A_0\).
CDF monotonicity throughout the middle interval implies
\[
 P_{\rm mid}\leq(1-2\kappa+u^2/2)\frac\ell2
 +\frac{\ell(4\kappa+2b_0+2b_1\log A_0+2uK)}{2A_0}
 +\frac{K^2\ell}{2A_0^2}.
\]

Start the discrete far sum at
\[
 L=\lfloor\sqrt{2M}\rfloor+1 .
\]
The initial interval \([\sqrt{2M},L]\) costs at most
\(1/(\sqrt2 A_0)\), including when \(\sqrt{2M}\) is an integer.
This prevents the false substitution \(f(Q)=0\) at \(T^2=2M\).
For all subsequent terms \(T^2>2M\).
Use \(\log(1+1/T)\leq1/T\), decreasing-sum comparison and
\(\sum_{T\geq L}T^{-2}\leq L^{-1}+L^{-2}\).
Equations \eqref{ap:far-integral} and \eqref{ap:deleted-error} then give
\[
 P_{\rm far}\leq
 \frac{1-\rho}{2}-\frac{3\alpha(1-\rho)}8
 +\frac{1+5\theta/4+4E(A_0)}{\sqrt2 A_0}
 +\frac{2E(A_0)}{A_0^2}.
\]
When \(Q<1\) the survival is zero; no negative logarithmic error is used.
Combining the three bounds proves
\begin{equation}
 \mu(M)\leq-\Delta(p,r)+\mathcal E_{p,r}(A_0),
 \label{ap:master-balance}
\end{equation}
where
\[
 \Delta=
 \rho\{\alpha(5/8+\ell)+1/2\}
 +3\alpha/8-1/2-\ell/2-\tfrac34(1-\ell)u^2
\]
and
\begin{align}
 \mathcal E_{p,r}(A_0)
 =\frac1{A_0}\bigl[&
 2\kappa+2(b_0-b_1)+2b_1\log A_0+2uK\beta \notag\\
 &+\tfrac\ell2(4\kappa+2b_0+2b_1\log A_0+2uK)
 +(1+5\theta/4+4E(A_0))/\sqrt2\bigr] \notag\\
 &+\frac{K^2\ell+2E(A_0)}{A_0^2}.
 \label{ap:master-error}
\end{align}
All these estimates hold for arbitrary positive prime valuations.
They are inequalities for the original finite average, not limiting-law
or monotonicity assertions.

\subsection{Uniform leading margins and universal cutoffs}
\label{ap:universal-cutoffs}
The margin \(\Delta\) increases with either prime.
To check this algebraically put \(x=1/p,y=1/r\),
\(A_*=\alpha(5/8+\ell)+1/2\), \(C_*=\tfrac34(1-\ell)\).
Multiplication of \(\partial_x\Delta\) by \((1+x)/\rho\) gives
\[
 -A_*+2C_*(1-x)(1+x)^2(1-y)^2(1+y)<0 .
\]
For \(0\leq x,y\leq1/2\), its positive term is \(<9/8\), whereas
\(A_*>51/40\), using \(\alpha>3/5\) and \(\ell>2/3\).
Thus both reciprocal-prime derivatives are negative.

The coarse coherent parameters \((b_0,b_1,K)=(4,3,4)\) and
\(E(A_0)=3+2\log A_0\) give
\begin{equation}
 \mathcal E_{p,r}(A_0)<
 \frac{22+15\log A_0}{A_0}+\frac{18+4\log A_0}{A_0^2}.
 \label{ap:coarse-error}
\end{equation}
For example, the coefficients of its linear constant and logarithm are
bounded by \(5167/240<22\) and \(141/10<15\); its quadratic constant
is \(<18\).
At the smallest corner,
\[
 \Delta(2,3)=\alpha(11/16+\ell/2)-1/3-5\ell/12
             >2885/521272>1/200.
\]
Indeed convexity bounds
\[
\zeta(2)<\sum_{n=1}^5n^{-2}+\int_{11/2}^{\infty}x^{-2}dx
=65159/39600.
\]
Insert \(\alpha>39600/65159\) first and
\(\ell<7/10\) second. Elementary finite positive exponential or atanh
series give
\[
 3/5<\alpha<2/3,\quad 2/3<\ell<7/10,\quad
 \log3<11/10,\quad\log5<17/10 .
\]
At \(A_0=2^{16}\), \(\log A_0<56/5\), and
\[
 \mathcal E<
 \frac{190}{65536}+\frac{314/5}{65536^2}<1/300.
\]
The positive error terms decrease thereafter.
Hence every two-prime modulus \(M\geq2^{32}\) has
\(\mu(M)<-1/600\).

The sharper three cutoffs follow from the refined parameters, still
without finite-height experiments:
\begin{equation}
 \begin{array}{ll}
 p=2,\ r\geq11,\ M\geq384^2:&\mu(M)<-1/600,\\
 p=3,\ r\geq7,\ M\geq512^2:&\mu(M)<-1/350,\\
 p\geq5,\ r\geq7,\ M\geq600^2:&\mu(M)<-1/750 .
 \end{array}
 \label{ap:three-cutoffs}
\end{equation}
Here are the exact parameter substitutions and endpoint comparisons.

For \(p=2\) the periodic extrema are \(h=2u,m=-u\);
they occur at \(n=r-2,r+1\). Thus \(K=4u,W=3u\).
For \(r\geq11\), using the second coherent primitive and density choices,
\[
 u\leq\tfrac12,\quad\kappa<\tfrac49,\quad\theta<\tfrac13,\quad
 b_0\leq\tfrac{11}4,\quad b_1\leq\tfrac34,\quad b_0-b_1\leq2,\quad K\leq2,
 \quad E_0\leq\tfrac{125}{44},\quad E_1\leq\tfrac34 .
\]
Substitution in \eqref{ap:master-error}, using
\(\beta<1/5\) and \(1/\sqrt2<3/4\), gives
\[
 \mathcal E<
 \frac{363/20+(171/40)\log A_0}{A_0}
 +\frac{17/2+(3/2)\log A_0}{A_0^2}.
\]
The unrounded linear constant is \(143521/7920<363/20\).
At \(A_0=384\), \(\log384<6\), so
\[
 \mathcal E<73/640+35/294912<23/200.
\]
The exact corner after inserting \(\alpha>3/5\), then \(\ell>2/3\), is
\[
 \Delta(2,11)>
 109/240-761/2178+(157/7260)(2/3)
 =3461/29040>7/60.
\]
This proves the first line of \eqref{ap:three-cutoffs}.

For \(p=3,r\geq7\), \eqref{ap:refined-discrepancy} gives
\[
 u\leq\tfrac23,\quad\kappa<\tfrac12,\quad\theta<\tfrac13,\quad
 K\leq\tfrac83,\quad b_0\leq\tfrac{14}3,\quad b_1\leq\tfrac43,\quad
 b_0-b_1\leq\tfrac{10}3,\quad E_0\leq\tfrac{250}{63},\quad E_1\leq\tfrac{10}9 .
\]
Therefore
\[
 \mathcal E<
 \frac{27+7\log A_0}{A_0}
 +\frac{13+(9/4)\log A_0}{A_0^2}.
\]
The unrounded linear constant is \(133843/5040<27\).
Since \(\log512<63/10\),
\[
 \mathcal E(512)<711/5120+1087/10485760<7/50.
\]
The corner computation is
\[
 \Delta(3,7)>
 603/1280-1307/3136+(1087/7840)(2/3)
 =27613/188160>1/7,
\]
which proves the second line.

Finally, at \(p=5,r=7\), insertion of \(\alpha>3/5\), then
\(\ell>2/3\), gives
\(\Delta(5,7)>95869/470400>61/300\).
The coefficient of \(\ell\) after the first insertion is positive.
Since \(600=2^3\cdot3\cdot5^2\), \(\log600<33/5\), and
\eqref{ap:coarse-error} gives
\[
 \mathcal E(600)<121/600+37/300000
                =60537/300000<101/500.
\]
The difference \(61/300-101/500=1/750\) proves the third line.
All endpoint comparisons are strict rational inequalities, and all error
terms decrease above the stated scale.

\paragraph{Coverage boundary.}
The three cutoffs cover every pair except radicals \(6,10,14,15\).
They reduce the remaining problem to a bounded domain, but do not themselves
certify that domain or those four families.
A complete proof of the all-modulus theorem must additionally supply their
separate family arguments, the later source-domain reductions and the exact
finite-certificate clauses. No sign extrapolation or monotonicity of the
true finite means may replace these remaining premises.


\section{The refined radical-six balance}\label{app:rad6}
This argument applies to every \(M\) with \(\rad(M)=6\), not just \(6^N\).
Put \(a=\sqrt M\), \(\alpha=\zeta(2)^{-1}\), \(\ell=\log2\),
\(\theta=\alpha/2\), and \(\theta_6=\alpha/12\).
Period six gives
\[
 -2/3\le U_6(x)-x/3\le2/3,\qquad
 -1/3\le U_6(n)-n/3\le2/3\quad(n\in\Z_{\ge0}).
\]
Partial summation yields
\(\sum_{d\le Q,(d,6)=1}1/d\le1+\log Q/3\).
For real \(Q\ge2\) the tail satisfies
\(Q\sum_{d>Q,(d,6)=1}d^{-2}\le1/3+4/(3Q)\le1\).
For \(1\le Q<2\), it is \(<2/9\) since
\(\sum_{d\ge2,(d,6)=1}d^{-2}=2\zeta(2)/3-1<1/9\).
The real-cutoff estimate is therefore valid without rounding \(Q\).

For \(x=n+t\), \(0\le t<1\),
\(|nU_6(n)-x^2/3|\le x\), by the displayed discrepancy.
Möbius inversion with only six-unit contents consequently gives
\[
 |P(T)-\alpha T^2/2|\le(4/3)T+(1/3)T\log T.        \tag{B.1}\label{eq:six-primitive}
\]
Every closed interval of length \(h\) contains at most \(h/3+4/3\)
six-units, including boundary atoms. One proof uses the alternating
spacings two and four: \(2m+1\) consecutive units span at least \(6m\)
and \(2m\) units span at least \(6m-4\).
For \(L_T=(2T-M/T)_+\), the collision triangle in the preceding appendix
then gives, when \(2T^2>M\),
\[
 J(T)\le L_T^2/18+(8/9)L_T+16/9.                 \tag{B.2}\label{eq:six-collision}
\]
At \(2T^2=M\), \(T\) is a nonunit, so the putative equality pair
\(d=f=T\) is forbidden. At \(T^2=M\), \(T\) is also a nonunit.
Thus the same endpoint conventions apply for unequal valuations.

There is a useful additional gain in the far tail. Define
\[
 A(Q)=\sum_{\substack{d\le Q\\(d,6)>1}}\varphi(d)/d,\qquad
 A_6(Q)=\sum_{\substack{d\le Q\\6\mid d}}\varphi(d)/d.
\]
Möbius expansion gives the exact identities
\[
 A(Q)=\sum_{\substack{d\le Q\\(d,6)=1}}\frac{\mu_{\rm Mob}(d)}d
 \left[\frac12\floor{Q/(2d)}+\frac23\floor{Q/(3d)}
                              -\frac56\floor{Q/(6d)}\right],
\]
\[
 A_6(Q)=\frac13\sum_{\substack{d\le Q/6\\(d,6)=1}}
             \frac{\mu_{\rm Mob}(d)}d\floor{Q/(6d)}.
\]
The first bracket minus \(Q/(3d)\) has period six and lies in
\([-2/3,1/3]\). The harmonic and real tail bounds above then give,
for real \(Q\ge1\),
\[
 |A(Q)-\theta Q|\le E(Q)=1+(2/9)\log Q,\quad
 |A_6(Q)-\theta_6 Q|\le E_6(Q)=2/3+(1/9)\log Q.    \tag{B.3}
\]
For example, decompose a denominator into its powers of two and three
and its six-unit part; summing the powers gives geometric series,
and the six-unit harmonic and tail estimates above bound the remaining
floor errors. Equivalently insert \(p=2,r=3\) into the exact deleted-density
floor identity in Appendix~\ref{ap:two-prime-compensation}, and separately
retain denominators divisible by both primes. The resulting coefficients
are \(\theta=\alpha/2\) and \(\theta_6=\alpha/12\).

For centers divisible by six the neighbor repair costs \(d\), rather
than \(2d\). With \(F\) defined in Section~\ref{sec:prime}, the main survival
envelope is therefore
\[
 V(t)=(\theta-\theta_6)F(t/\sqrt2)+2\theta_6F(t)
                                      +\theta/(2t^2),\quad t\ge\sqrt2 .
\]
In partial summation, the difference between the two rational kernels
is unimodal; its maximum is at most \(2M^2/T^3\) and its total variation
is at most twice this. The extra error is consequently
\(8E_6(M/T)M/T^3\); the original one is \(4E(M/T)/T\).
Termwise positive-series integration, or the identity
\(2\Li(1/2)=\zeta(2)-\ell^2\), gives
\[
 \int_{\sqrt2}^{\infty}V(t)\frac{dt}{t}
       =1/4-3\alpha/16-\alpha(1-\ell)^2/24 .
\]
The dilogarithm identity follows by differentiating
\(\Li(x)+\Li(1-x)=\zeta(2)-\log x\log(1-x)\) and fixing the endpoint.

Repeat the central and middle integrals of the preceding appendix with
\eqref{eq:six-primitive}--\eqref{eq:six-collision}, and start the far sum
at \(\floor{\sqrt{2M}}+1\).
The complete result is
\[
 \mu(M)\le-\Delta_6+\mathcal E_6(a),
\]
where
\[
 \Delta_6=\alpha(11/16+\ell/2)-1/3-5\ell/12
                                +\alpha(1-\ell)^2/24>1/130
\]
and, writing \(\beta=3-2\sqrt2\),
\begin{align*}
 \mathcal E_6(a)=\frac1a\bigl[&
 \alpha+2+(2/3)\log a+(8/9)\beta\\
 &+\ell(2\alpha+32/9+(2/3)\log a)/2\\
 &+(1+5\alpha/8+4E(a)+(4/3)E_6(a))/\sqrt2\bigr]\\
 &+\{16\ell/9+2E(a)+2E_6(a)\}/a^2 .
\end{align*}
The margin follows from
\(\alpha>39600/65159,\ \ell<7/10\):
the part without the final gain exceeds \(2885/521272\),
and the gain exceeds \(1/440\), whose sum exceeds \(1/130\).
At \(a_0=6^{9/2}>3000\), \(\log a_0<81/10\);
using \(\alpha<2/3,\ell<7/10,1/\sqrt2<3/4,\beta<1/5\) gives
\[
 \mathcal E_6(a_0)<23/a_0+10/a_0^2<6901/900000<1/130.
\]
Every error term decreases for \(a\ge a_0\). Hence
\(\mu(M)<-287/11700000\) whenever \(\rad(M)=6,\ M\ge6^9\).
The eight smaller diagonal moduli \(6^1,\ldots,6^8\) are finite
certificates; unequal valuations below the cutoff are handled separately.

\section{Elementary constant enclosures}\label{app:constants}
All nonrational constants used for cutoff comparisons admit positive
series enclosures. In particular
\[
 \sum_{n=1}^{10}n^{-2}+\frac1{11}>\frac{41}{25},\qquad
 \zeta(2)<\sum_{n=1}^{5}n^{-2}+\frac2{11}
          =\frac{65159}{39600}<\frac{33}{20}.
\]
For the upper bound, convexity of \(x^{-2}\) bounds the term at each
integer \(n\ge6\) by its integral on \([n-1/2,n+1/2]\).
The lower bound uses the integral tail after ten terms.
Further,
\[
 e\le\sum_{n=0}^{6}\frac1{n!}
             +\frac{1}{7!(1-1/8)}<\frac{87}{32}<3.
\]
The atanh enclosure in Section~\ref{sec:cert}, applied to rational
arguments, proves all logarithmic comparisons with outward rounding.
For instance its range reduction gives
\(2/3<\log2<7/10,\ \log3<11/10,\ \log5<17/10\),
\(\log140<5,\ \log384<6,\ \log512<63/10,\ \log600<13/2\).
Alternatively, lower partial sums of \(\exp x\) prove each upper
logarithm comparison. These checks do not rely on floating decimal
evaluation.


\section{The short-length bridges}\label{app:shortlength}
Here we make explicit the predecessor result at lengths at least two
used in Appendix~\ref{ap:finite-two-prime}. This order avoids a circular
use of the length-one completion.
First the three sixth-power certificates for \(B=10,12,14\), the eight
base-six certificates, and the coarse cutoffs give every two-prime base
at \(N\ge6\), as explained in the family proof.
The complete radical families then follow without using any
length-two or length-one result.

Outside radicals \(6,10,14,15\), the coarse balance admits further
corner cutoffs. Monotonicity of its leading margin gives
\[
 \begin{gathered}
 \Delta(2,11)>3461/29040>1/9,\\
 \Delta(3,7)>27613/188160>1/7,\quad
 \Delta(2,17)>1447/10404>1/8 .
 \end{gathered}
\]
At \(a=1296=6^4\), using \(\log a<36/5\), the coarse error
\((22+15\log a)/a+(18+4\log a)/a^2\)
is \(<46813/466560<1/9\).
At \(a=1024\), using \(\log a<7\), it is
\(<65047/524288<1/7\) and \(<1/8\).
These comparisons and decreasing errors give:
\[
 \begin{array}{c|c}
 \hbox{remaining prime factors}&\hbox{sufficient }M\\\hline
 p=2,\ r=11\hbox{ or }13&6^8\\
 p=2,\ r\ge17&2^{20}\\
 p\ge3,\ r\ge7&2^{20}.
 \end{array}                                                    \tag{D.1}
\]
In particular all lengths \(N\ge5\) are covered.
For length four the complete remaining bases are \(21,22,26\);
their three finite certificates prove all lengths \(N\ge4\).

For length three define the conservative domain of bases \(B=p^a r^b\)
outside the four complete radicals by
\(B^3<6^8\) if \(p=2,r=11,13\), and \(B^3<2^{20}\) otherwise.
Remove bases already covered by reindexing at length at least four.
This gives exactly the following 35 bases:
\[
 \begin{gathered}
 21,22,26,33,34,35,38,39,44,46,51,52,55,57,58,62,63,65,68,69,\\
 74,76,77,82,85,86,87,88,91,92,93,94,95,99,104 .
 \end{gathered}
\]
The retained domain generator also gives their unique prime valuations.
Every one has a negative certificate, proving all lengths \(N\ge3\).
The numerical limits behind the enumeration are \(B\le118\) in the
first category and \(B\le101\) in the others.

For length two enumerate every integer base \(B\le1295\) with exactly
two prime factors outside the four radicals. Keep
\[
 B<1296\quad(p=2,r=11,13),\qquad
 B<1024\quad\hbox{otherwise},
\]
and remove \(B=441,484,676\), which are respectively the squares of
the already certified bases \(21,22,26\). Any other higher-power
reindexing in this range is either in a complete radical family or
above its analytic cutoff.
There are exactly 454 retained bases, split into 256 odd bases,
183 bases with \(p=2,r\ge17\), and 15 with \(p=2,r=11,13\).
Their physical modulus is \(M=B^2\); they are not 454 arbitrary
two-prime moduli.
The parameterized balance with the actual periodic discrepancy and
outward constant enclosures pays 355 of them. The other 99 have
negative histogram/log certificates constructed by Section~\ref{sec:cert}.
This proves every two-prime base at length two.
All lengths \(N\ge2\) are now proved before the length-one residual
inventory is formed.

The arithmetic in the analytic filters is also fully finite.
For \(\zeta(2)\) use
\[
 \sum_{j=1}^{K}j^{-2}+1/(K+1)<\zeta(2)
            <\sum_{j=1}^{K}j^{-2}+1/K ,
\]
invert outward to enclose \(\alpha\), and enclose logarithms by the
atanh series. Square roots have rational enclosures obtained by integer
square roots. The length-two filter uses \(K=128\) for the zeta sum;
the companion records the root scale \(10^{15}\), logarithmic precision
and any outward dyadic grid separately. A row is paid only if its entire upper enclosure is
strictly negative; an inconclusive enclosure remains in the finite
certificate set. The later length-one five-bound filter is evaluated
in the same outward manner. A rounded decimal sign is never a filter.

\begin{table}[ht]
\centering
\begin{tabular}{lrr}\toprule
Bridge & Rows & Kind of retained evidence\\\midrule
Odd prime powers &21&exact heights and integer Jensen sums\\
Base six, \(N=1,\ldots,8\)&8&upper histograms and integer log sums\\
Length six, bases \(10,12,14\)&3&upper histograms and integer log sums\\
Squarefree semiprime prefix&4596&digests/signs; reconstruction code\\
Offdiagonal radical six&143&upper histograms and log ledger\\
Radical ten&97&upper histograms and log ledger\\
Radicals fourteen/fifteen&132&upper histograms and log ledger\\
Length four&3&upper histograms and log ledger\\
Length three&35&upper histograms and log ledger\\
Length two residual&99&upper histograms and log ledger\\
Final length one&4410&upper histograms and log ledger\\\bottomrule
\end{tabular}
\caption{Finite proof dependencies. Domains are specified above and in the
companion. Rows with different proof roles may overlap and are not to be
added as a count of distinct moduli.}\label{tab:dependencies}
\end{table}


\section{Effective reduction and finite completion for two prime factors}
\label{ap:finite-two-prime}

Write
\[
 M=p^a q^b,\qquad p<q\ \hbox{distinct primes},\qquad a,b\ge1,
 \qquad
 \mu(M)=\frac1M\sum_{s=0}^{M-1}\log R_M(s)-\frac12\log M .
\]
The minimum defining \(R_M\) uses positive denominators coprime to the
whole \(M\), and zero is one of the \(M\) equally weighted states.
All logarithms in this appendix are natural.
The reductions below apply to the actual valuations \(a,b\); they do not
compare means on different moduli.
We write \(\mu_{\mathrm{Mob}}\) for the M\"obius function, to distinguish
it from the centered mean \(\mu(M)\).

\subsection{Exact discrepancy parameters and admissible errors}
\label{ap:finite:parameters}

Put
\[
 \begin{gathered}
 D=pq,\quad u=\frac{(p-1)(q-1)}{pq},\quad
 \rho=\frac{pq}{(p+1)(q+1)},\\
 \alpha=\frac1{\zeta(2)},\quad \ell=\log2,\quad
 \kappa=\alpha\rho,\quad \theta=\alpha(1-\rho),\quad
 \beta=3-2\sqrt2 .
 \end{gathered}
\]
For integer \(n\ge0\), let
\[
 U_D(n)=n-\lfloor n/p\rfloor-\lfloor n/q\rfloor+\lfloor n/D\rfloor,
 \qquad e(n)=U_D(n)-un .
\]
The function \(e\) is \(D\)-periodic and
\[
 e(n)-e(n-1)=\mathbf1_{(n,D)=1}-u,\qquad
 e(D-n)=\mathbf1_{(n,D)=1}-e(n)\quad(0<n<D).
\]
In particular,
\begin{equation}
 -u\le e(n)\le2u. \label{eq:finite:discrepancy}
\end{equation}
Indeed a minimum occurs at zero or a nonunit.  For \(n=pi<D\),
write \(pi=tq+r\), \(1\le r<q\), \(0\le t<p\).  Then
\[
 e(pi)=\frac{r-i}{q},\qquad
 i-r=tq-(p-1)i\le\frac{tq-(p-1)}p
                  \le\frac{(p-1)(q-1)}p=uq .
\]
The other nonunit family is treated with the primes interchanged.
Reflection bounds a nonunit above by \(u\).  For a unit, its reflected
unit follows a step of size \(1-u\) from a value at least \(-u\);
reflection then bounds the original value above by \(2u\).

Let \(\sigma\) be the maximum of \(e(n)\) over nonunits in \(0\le n<D\),
including \(n=0\).  The minimum over the whole period is \(-\sigma\),
and its maximum is \(u+\sigma\): a maximum is immediately before a
nonunit step, including the wraparound predecessor \(D-1\).
An exact short formula, not a scan of all \(D\) residues, is
\begin{align}
 i_t&=\min\!\left(q-1,\left\lfloor\frac{q(t+1)-1}{p}\right\rfloor\right),
       &&0\le t<p, \notag\\
 \sigma&=\max\left(
 0,\ \max_{0\le t<p}\frac{(p-1)i_t-qt}{q},\
       \max_{1\le j<p}\frac{(qj\bmod p)-j}{p}\right).
 \label{eq:finite:sigma}
\end{align}
To prove this, the nonunits are \(pi\), \(1\le i<q\), and \(qj\),
\(1\le j<p\).  On the block \(t=\lfloor pi/q\rfloor\),
the first numerator is increasing in \(i\), so its maximum is at \(i_t\).
Every such block is nonempty since \(q>p\); the minus one in
\eqref{eq:finite:sigma} preserves its strict upper endpoint.
There are at most \(2p-1\) nonzero candidates.

Define
\[
 K=2u+2\sigma,\qquad W=2u+\sigma .
\]
The real discrepancy lies between \(-\sigma-u\) and \(u+\sigma\).
Consequently a closed interval of length \(h\) contains at most \(uh+K\)
positive \(D\)-units, with both boundary atoms included.  Also
\[
 \sum_{\substack{d\le Q\\(d,D)=1}}\frac1d
       \le u+(u+\sigma)+u\log Q\qquad(Q\ge1).
\]
For \(x=n+t\), \(n=\lfloor x\rfloor\), \(0\le t<1\), the same discrepancy
bounds give \(\left|nU_D(n)-ux^2\right|\le Wx\).
M\"obius inversion and its integer-cutoff tail therefore yield
\begin{equation}
 \left|P(T)-\kappa T^2\right|\le T(b_0+b_1\log T),\qquad
 P(T)=\sum_{\substack{d\le T\\(d,D)=1}}\mu_{\mathrm{Mob}}(d)
       \lfloor T/d\rfloor U_D(\lfloor T/d\rfloor),
 \label{eq:finite:primitive-error}
\end{equation}
with either admissible pair
\begin{align}
 (b_0,b_1)&=(3u+\sigma,\,2u+\sigma), \label{eq:finite:pair-one}\\
 (b_0,b_1)&=((2u+\sigma)^2+u,\,u(2u+\sigma)).
 \label{eq:finite:pair-two}
\end{align}
The coefficient follows from
\[
 u\sum_{(d,D)=1}\frac{\mu_{\mathrm{Mob}}(d)}{d^2}
 =\frac{u\alpha}{(1-p^{-2})(1-q^{-2})}=\kappa .
\]
For the first pair use \(\sum_{d\le T}d^{-1}\le1+\log T\);
for the second use the displayed unit harmonic bound.  The omitted series
costs at most \(uT\).  Thus \(b_0-b_1\) is respectively \(u\) and
\((2u+\sigma)(u+\sigma)+u\), both nonnegative.
The independent coarse choice \((b_0,b_1,K)=(4,3,4)\) is also valid.

For deleted denominators, set
\[
 A(Q)=\sum_{\substack{h\le Q\\(h,D)>1}}\frac{\varphi(h)}h,\quad
 v=2-\frac1p-\frac1q,\quad
 c_0=\frac{p-1}{p^2}+\frac{q-1}{q^2}-\frac{D-1}{D^2}.
\]
Decomposing squarefree totient divisors gives the exact identity
\[
 A(Q)=\sum_{\substack{d\le Q\\(d,D)=1}}\frac{\mu_{\mathrm{Mob}}(d)}d
 \left[
 (1-1/p)\left\lfloor\frac Q{pd}\right\rfloor+
 (1-1/q)\left\lfloor\frac Q{qd}\right\rfloor-
 (1-1/D)\left\lfloor\frac Q{Dd}\right\rfloor
 \right].
\]
The full sum's leading coefficient is \(\theta\).
The bracket's linear coefficient is \(c_0Q/d\), and its error is at most \(v\);
for a real cutoff, \(Q\sum_{d>Q}d^{-2}<2\), not merely one.
Hence either of the complete errors
\begin{equation}
 E_{\rm r}(x)=v(2u+\sigma)+2c_0+vu\log x,
 \qquad E_{\rm u}(x)=3+2\log x
 \label{eq:finite:deleted-errors}
\end{equation}
pays \(|A(Q)-\theta Q|\) for real \(Q\ge1\).
The refined constant also bounds the below-one partial-summation interval:
\(\theta<1/3\), whereas \(vu\ge7/18\).
In particular,
\[
 \sum_{\substack{h\le Q\\(h,D)>1}}\varphi(h)
       \le \frac{\theta Q^2}{2}+2QE(Q).
\]

\subsection{Finite balance, closed endpoints and the four simultaneous gains}
\label{ap:finite:five-bounds}

For every integer \(T^2\le M\), the original central count satisfies
\[
 C_M(T)=1+2P(T)-J_M(T).
\]
Reduction of common content is legitimate because it divides a full-unit
denominator.  Same-sign primitive images are injective.  Opposite-sign
images meet precisely at \(Ad+Eb=M\), with positive \(A,E,b,d\le T\).
Thus \(b+d\ge M/T\), with equality retained, and \((b,d)=1\).
For a fixed ordered denominator pair the allowed phase interval has width
at most \([T(b+d)-M]/(bd)\le1\); equality would require \(M=T^2\) and
a denominator \(T\), which is nonunit at a square modulus.
Putting \(L_T=(2T-M/T)_+\), the closed interval unit bound gives
\[
 J_M(T)\le\frac{u^2L_T^2}{2}+2uKL_T+K^2\quad(2T^2>M).
\]
If \(2T^2=M\), equality would force \(b=d=T\); the odd prime divisor
of \(M\) divides \(T\), so the collision is absent.  Thus \(J_M(T)=0\)
also at that boundary.

Here is the far-cover argument used in the finite balance.
For \(T^2>2M\), full Farey neighbors of order \(T\) have coprime
denominators and sum greater than \(T\).  Two retained neighbors cover
their intervening arc.  If both are deleted, their prime obstructions
are complementary, so their denominator difference is a full unit;
when both denominators are at least \(M/T\), that difference witness
covers the arc.  An uncovered point must therefore touch a deleted
center \(c/h\) with \(h<M/T<T/2\).
Its complete neighbors have denominators greater than \(T-h\);
subtracting \(c/h\) once repairs a bad neighbor and gives a determinant-one
unit neighbor of denominator greater than \(T-2h\).  When \(D\mid h\),
the complete neighbors were already units.  The corresponding open
hole has length at most
\[
 \frac{2(M-hT)}{h(T-j_hh)},\qquad
 j_h=\begin{cases}1,&D\mid h,\\2,&D\nmid h.\end{cases}
\]
One additional integer per reduced center pays rounding; overlaps only
enlarge the upper bound.  Therefore
\begin{equation}
 M-C_M(T)\le
 \sum_{\substack{2\le h<M/T\\(h,D)>1}}\varphi(h)+
 \sum_{\substack{2\le h<M/T\\(h,D)>1}}
       \frac{2\varphi(h)(M-hT)}{h(T-j_hh)} .
 \label{eq:finite:far-cover}
\end{equation}
The denominators are strictly positive.  No center phase or exact tail
distribution is needed.

Write \(x=\sqrt M\), \(k=\lfloor x\rfloor\), \(t=x-k\).
For any compatible primitive pair and deleted error \(E\), define
\begin{align}
 \Delta&=\rho\{\alpha(5/8+\ell)+1/2\}
          +3\alpha/8-1/2-\ell/2-\tfrac34(1-\ell)u^2,
 \label{eq:finite:delta}\\
 {\cal E}_{b,E}(x)
 &=\frac{2\kappa+2(b_0-b_1)+2b_1\log x+2uK\beta}{x}\notag\\
 &\quad+\frac\ell{2x}(4\kappa+2b_0+2b_1\log x+2uK)\notag\\
 &\quad+\frac{1+5\theta/4+4E(x)}{\sqrt2 x}
 +\frac{K^2\ell+2E(x)}{M}. \label{eq:finite:error}\\
 \delta_x&=\frac{1+2t-3t^2}{2}+\frac{t^3}{x},\qquad
 G_x=2x(1-t)+t^2=2x-M+k^2, \notag\\
 {\cal G}_b(x)&=
 \frac{\ell\kappa G_x+2\kappa\delta_x
                  +2(b_0-b_1)+\log x}{M}. \label{eq:finite:gain}
\end{align}
Then the strengthened finite balance is
\begin{equation}
 \mu(M)\le-\Delta+{\cal E}_{b,E}(x)-{\cal G}_b(x).
 \label{eq:finite:strengthened}
\end{equation}

For completeness, its integration and rounding costs can be checked directly.
The centered mean is minus the central CDF integral up to \(x\), plus
the survival integrals over \([x,\sqrt2x]\) and \([\sqrt2x,\infty)\).
The collision integrals are
\[
 \frac1M\int_{x/\sqrt2}^x(2y-M/y)^2\frac{dy}{y}
       =\frac32-2\ell,\qquad
 \int_{x/\sqrt2}^x(2y-M/y)\frac{dy}{y}=x\beta .
\]
The central credit and monotone middle plateau, before retaining gains,
are respectively at least
\[
 \kappa-u^2(3/4-\ell)
 -\frac{2\kappa+2(b_0-b_1)+2b_1\log x+2uK\beta}{x}
 -\frac{K^2\ell}{2M}
\]
and at most
\[
 (1-2\kappa+u^2/2)\frac\ell2
 +\frac{\ell(4\kappa+2b_0+2b_1\log x+2uK)}{2x}
 +\frac{K^2\ell}{2M}.
\]
For the far sum use \(h<M/T\), partial summation of \(A\), and
\[
 F(z)=\sum_{n\ge1}\frac{z^{-2n}}{n(n+1)},\qquad F(1)=1.
\]
The leading normalized survival is bounded by
\(\theta F(T/(\sqrt2x))+\theta x^2/(2T^2)\).
Its logarithmic integral is \((1-\rho)/2-3\theta/8\).
The first integer far threshold is \(\lfloor\sqrt{2M}\rfloor+1\);
the initial interval costs at most \(1/(\sqrt2x)\), including equality
at an integer \(\sqrt{2M}\).  The first summand and discrete error add
\((5\theta/4+4E(x))/(\sqrt2x)+2E(x)/M\).
These terms give \eqref{eq:finite:delta}--\eqref{eq:finite:error}.

The four gains in \eqref{eq:finite:gain} affect distinct terms.
Indeed
\[
 \int_1^x\lfloor y\rfloor^2\frac{dy}{y}
 \ge\frac{(k-1)^2}{2}+\frac{k^2t}{x}
   =\frac{x^2}{2}-x+\delta_x .
\]
Here \(\log(1+1/j)\ge1/j-1/(2j^2)\) pays the complete intervals and
\(\log(x/k)\ge t/x\) pays the last interval.
Exact primitive-error integration restores \(2(b_0-b_1)/M\);
the isolated zero restores \((\log x)/M\).
Retaining the exact endpoint \(k^2\), instead of \(M-2x\), restores
\(\ell\kappa G_x/M\) on the middle plateau.
Collision and far penalties are unchanged, so all four improvements
coexist.  At a square, \(\delta_x=1/2\), \(G_x=2x\).
They must remain paired with their own \(b_0,b_1\).

For the coarse pair \((4,3,4)\), the error in \eqref{eq:finite:error} is
strictly below
\[
 {\cal E}_{\rm c}(x)=\frac{22+15\log x}{x}
                           +\frac{18+4\log x}{M}.
\]
This follows from
\(2\kappa\le4/3\), \(\theta\le1/3\), \(u\le1\),
\(\ell<7/10\), \(\beta<1/5\), \(1/\sqrt2<3/4\):
the linear constant is at most \(5167/240<22\), its logarithmic
coefficient is at most \(141/10<15\), and the quadratic constant is
less than \(18\).
Thus the following are five separately valid, complete upper bounds:
\[
 -\Delta+{\cal E}_{\rm c}-{\cal G}_{(4,3)},\qquad
 -\Delta+{\cal E}_{b,E}-{\cal G}_b,\quad
 b\in\{\eqref{eq:finite:pair-one},\eqref{eq:finite:pair-two}\},\
 E\in\{E_{\rm r},E_{\rm u}\}.
\]
Taking their minimum is legitimate; moving a gain from one bound to
another is not.

\subsection{Uniform modulus cutoffs}
\label{ap:finite:cutoffs}

The leading credit increases with either prime parameter.  Put
\(z=1/p,w=1/q\), \(A_0=\alpha(5/8+\ell)+1/2\),
\(C_0=3(1-\ell)/4\).  Multiplying its \(z\)-derivative by
\((1+z)/\rho>0\) gives
\[
 -A_0+2C_0(1-z)(1+z)^2(1-w)^2(1+w)<0.
\]
Indeed the two factors are at most \(9/4\) and \(1\), respectively,
\(2C_0<1/2\), and \(A_0>51/40>9/8\), using
\(\alpha>3/5\) and \(2/3<\ell<7/10\).
This is monotonicity of the bound, not of the true mean.

Since \(\sigma\le u\), the second primitive pair has
\[
 b_0\le9u^2+u,\quad b_1\le3u^2,\quad
 b_0-b_1\le6u^2+u,\quad K\le4u .
\]
For \(p=2\), reflection gives the exact values
\(\sigma=u\), \(K=4u\), \(W=3u\).
Using the refined deleted error in \eqref{eq:finite:error}, elementary
substitution gives the following bounds and rational endpoint witnesses:
\[
\begin{array}{c|c|c|c}
\hbox{prime range}&\hbox{error upper}&x_0&
                   \hbox{leading credit lower}\\ \hline
p=2,\ q\ge11&
\dfrac{363/20+(171/40)\log x}{x}
 +\dfrac{17/2+(3/2)\log x}{x^2}&384&
\dfrac{3461}{29040}>\dfrac7{60}\\[5pt]
p=3,\ q\ge7&
\dfrac{27+7\log x}{x}
 +\dfrac{13+(9/4)\log x}{x^2}&512&
\dfrac{27613}{188160}>\dfrac17\\[5pt]
p\ge5,\ q\ge7&{\cal E}_{\rm c}(x)&600&
\dfrac{95869}{470400}>\dfrac{61}{300}
\end{array}
\]
For the first two rows use respectively
\[
 c_0\le13/44,\ E_{{\rm r},0}\le125/44,\ E_{{\rm r},1}\le3/4,
 \quad
 c_0\le20/63,\ E_{{\rm r},0}\le250/63,\
 E_{{\rm r},1}\le10/9 .
\]
For the second row \(u\le2/3\), \(\kappa<1/2\),
\(K\le8/3\), \(b_0\le14/3\), \(b_1\le4/3\),
\(b_0-b_1\le10/3\).  The first row uses
\(u\le1/2\), \(\kappa<4/9\), \(K\le2\),
\(b_0\le11/4\), \(b_1\le3/4\), \(b_0-b_1\le2\).
The resulting unrounded linear constants are
\(143521/7920<363/20\) and \(133843/5040<27\).
At the three endpoints,
\[
\begin{split}
 {\cal E}(384)&<73/640+35/294912<23/200,\\
 {\cal E}(512)&<711/5120+1087/10485760<7/50,\\
 {\cal E}_{\rm c}(600)&<60537/300000<101/500.
\end{split}
\]
Use \(\log384<6\), \(\log512<63/10\), \(\log600<33/5\);
these follow from the positive exponential sums giving
\(\log2<7/10,\log3<11/10,\log5<17/10\).
The constant and logarithmic error terms decrease thereafter.
It follows, for arbitrary positive valuations, that
\begin{equation}
 \begin{array}{ll}
 p=2,\ q\ge11,\ M\ge147456:&\mu(M)<-1/600,\\
 p=3,\ q\ge7,\ M\ge262144:&\mu(M)<-1/350,\\
 p\ge5,\ q\ge7,\ M\ge360000:&\mu(M)<-1/750.
 \end{array} \label{eq:finite:cutoffs}
\end{equation}
Equality at these sufficient cutoffs is already covered.

\subsection{The larger prime to the first power}
\label{ap:finite:ratio}

Suppose \(M=Kq\), \(K=p^a\), \(q>K\), \(M\ge19600\).
In the single-\(p\)-mark model, all central denominators are smaller
than \(q\) and are therefore also full units.  Primitive cancellation
and the collision-free range through \(\lfloor\sqrt{M/2}\rfloor\)
transfer unchanged from the single-mark proof.  In its Farey tail,
adjacent denominators cannot both be multiples of \(p\); its deleted
center bracket has denominators greater than \(T-h\).
Thus the odd-prime estimates of Section~\ref{sec:prime}, and the binary
calculation given in the following squarefree subsection, transfer to
this modulus without a prime-power assumption:
\[
\begin{array}{c|c|c}
p&\hbox{leading credit}&\hbox{error at }x=\sqrt M\\ \hline
2&1/12&(33/8+(12/11)\log x)/x+1/(4x^2)\\
3&41/486&(107/27+(16/11)\log x)/x+1/(3x^2)\\
\ge5&91/800&(151/40+(24/11)\log x)/x+1/(2x^2)
\end{array}
\]
At \(x=140\), \(\log140<5\) bounds the errors by quantities strictly
below \(7/100,82/1000,106/1000\).  Explicit upper fractions are
\[
 59021/862400,\quad467279/5821200,\quad22619/215600 .
\]
All three errors decrease for \(x\ge140\).  Their strict single-mark
margins are respectively greater than \(1/75,287/121500,31/4000\).

Below real height \(q\), single-mark and full-unit events coincide.
For a complete reduced lattice basis, the two-prime short-basis argument
gives \(R_M\le\sqrt2M/\lambda_1\).  The all-state incidence bound
\(\Pr(\lambda_1\le Y)\le3Y^2/M\) therefore implies
\(\Pr(R_M>X)\le\min(1,6M/X^2)\).
If \(q\ge6K\), the extra full-unit logarithmic cost is at most
\[
 \int_q^\infty\frac{6M}{X^2}\frac{dX}{X}=\frac{3K}{q}.
\]
Consequently
\begin{equation}
 b=1,\quad M\ge19600,\quad q\ge c(p)p^a
 \quad\Longrightarrow\quad\mu(M)<0,\qquad
 c(p)=\begin{cases}225,&p=2,\\1300,&p=3,\\400,&p\ge5.\end{cases}
 \label{eq:finite:ratio}
\end{equation}
These ratios imply \(q>K\) and \(q\ge6K\).
The comparisons are \(3/225=1/75\),
\(287/121500-3/1300=43/789750>0\), and
\(31/4000-3/400=1/4000>0\).
The first equality still pays a strict result because its single-mark
margin is strict.  The proof has \(b=1\); it is not a radical-only rule
and cannot be applied with the prime roles reversed.

\subsection{Squarefree semiprimes: zero mark loss and the complete prefix}
\label{ap:finite:squarefree}

Let \(R_{p;M}\) denote the minimum with only \(p\nmid d\).
Separating witnesses according to whether \(q\mid d\) gives exactly
\[
 R_{p;M}(s)=\min\{R_M(s),\
                 qR_{p;M/q}(s\bmod M/q)\}.
\]
If \(q\mid d\), the congruence forces \(q\mid A\), and division by \(q\)
gives the displayed child congruence; conversely multiplying a child
witness by \(q\) gives a sole-\(p\) witness.  No arbitrary nonunit
content is cancelled.  At \(M=pq\), the circle-gap cap proved below
is \(R_{pq}\le q\), while the second term is at least \(q\).
Consequently \(R_{pq}=R_{p;pq}\) pointwise.
Since \(\sqrt{pq}<q\), the single-mark transfer of the preceding
subsection applies, with NO mark-loss penalty.  Its odd-prime
estimates give \(\mu(pq)<-1/500\) for \(p\ge3\), \(pq\ge19600\).

For the even exterior we spell out the smaller sufficient threshold.
The binary single-mark calculation below is valid more generally for
\(M=2^a q\), \(q>2^a\), with \(x=\sqrt M<q\): all its central
denominators are full units, and the sole-two-mark Farey tail depends
only on the excluded even denominators. In particular it applies at
\(M=2q\). The divisor union bound on pairs with odd denominator gives
\[
 P_2(T)\ge
 [1-3\zeta(2)/8]T^2-\tfrac12T\log T
 \ge\tfrac38T^2-\tfrac12T\log T.
\]
Indeed there are at least \(T^2/2\) raw pairs; remove every odd
content \(h\ge3\), bounding its count by
\(T^2/(2h^2)+T/(2h)\).
The sums are bounded by
\(\sum_{h\ge3,\ h\ {\rm odd}}h^{-2}=3\zeta(2)/4-1\) and
\(\sum_{3\le h\le T,\ h\ {\rm odd}}h^{-1}\le\log T\).
Here \(\zeta(2)<5/3\).
For \(T\le K=\lfloor x/\sqrt2\rfloor\), \(x=\sqrt M\), there are no
collisions, and hence \(C_M(T)\ge cT^2-dT\log T\) with \(c=3/4,d=1\).
The quadratic floor integral and monotone plateau give
\[
 \frac1M\int_1^x C_M(\lfloor y\rfloor)\frac{dy}{y}
 \ge c\frac{K^2}{x^2}\left(\frac12+\log\frac{x}{K}\right)
 -\frac{2cK}{x^2}
 -d\frac{K\log K}{x^2}\left(1+\log\frac{x}{K}\right).
\]
This follows by integrating to \(K\), using
\(\int_1^K\lfloor y\rfloor^2dy/y\ge K^2/2-K\), bounding the
primitive-error integral by \(K\log K\), and retaining the plateau
\(C_M(K)\log(x/K)\).  The displayed \(2cK\) is a safe enlargement
of the floor loss.
For \(x\ge115\),
\[
 \begin{gathered}
 K/x>7/10-1/115>11/16,\quad K/x<3/4,\\
 K^2/x^2\ge1/2-\sqrt2/x,\quad
 \ell/2\le\log(x/K)\le5/11 .
 \end{gathered}
\]
With \(2/3<\ell<1\) and \(\sqrt2<3/2\), this yields central credit
\(5/16-[9/4+(12/11)\log x]/x\).
The single-mark Farey tail has majorant
\[
 V(z)=\tfrac12F(z)+\frac1{8z^2},\quad z\ge1,\qquad
 V(1)=5/8,\qquad
 \int_1^\infty V(z)\frac{dz}{z}
 =\frac{\zeta(2)-1}{4}+\frac1{16}<\frac{11}{48}.
\]
The initial interval costs \(1/x\), the first summand \(5/(8x)\),
and the integer-hole rounding costs \(1/(4x)+1/(4x^2)\).
Subtracting the central credit gives, for every \(x\ge115\),
\[
 \mu(2q)\le-\frac1{12}
 +\frac{33/8+(12/11)\log x}{x}+\frac1{4x^2}.
\]
The error decreases there.  Since \(\log115<5\), its remaining
strict margin at \(x=115\) is
\[
 \frac1{12}-\frac{843}{10120}-\frac1{52900}
       =\frac{1127}{80302200}>0 .
\]
Thus every odd prime \(q\ge6613\) lies in this analytic exterior:
\(2q\ge13226>115^2\).

The complete prefix is exactly
\[
 \{2q:3\le q<6613,\ q\ {\rm prime}\}
 \ \cup\
 \{pq:3\le p<q,\ p,q\ {\rm prime},\ pq<19600\}.
\]
It has \(853\) even and \(3743\) odd moduli, \(4596\) in total.
Every row has a negative outward-log certificate obtained from a
legal full-unit upper histogram, capped at the proved \(q\);
the pointwise construction is proved in Section~\ref{sec:early-histograms},
and the logarithmic certificate inequality is given below.
This finite prefix plus the analytic exterior proves every
squarefree-semiprime mean strictly negative, without any
transfer to higher valuations or any sampling argument.

\subsection{Complete small-radical families and their finite witnesses}
\label{ap:finite:families}

For clarity, the earlier complete families have an explicit infinite
part and a separately retained finite part.  The uniform coarse balance
gives the cutoffs \(2^{22}\) for odd prime pairs and \(9\,000\,000\)
for \(p=2,q\ge5\): at \(x=2048\) its error is less than \(7/100\) and
\(\Delta(3,5)>1691/14400>1/10\); at \(x=3000\) its error is less than
\(1/20-2026/937500\) and \(\Delta(2,5)>263/4500>1/20\).

The separately proved refined radical-six finite balance gives
\(\mu(M)<-287/11700000\) for every modulus of radical six with
\(M\ge6^9\).  It retains the shorter repaired bracket at centers
divisible by both primes, and applies to unequal valuations, with
nonunit central equality endpoints and a strictly started far sum.
This preceding radical-six lemma, rather than an asymptotic sign or
the coarse error above, is used for that one infinite family.

The following finite witness table completes the families.
Each entry means that every original modulus in the specified set has
a retained negative outward-log upper certificate, not that its true
minimum distribution was fully enumerated.
\[
\begin{array}{c|l|r}
\hbox{family}&\hbox{finite set below its analytic cutoff}&\hbox{rows}\\ \hline
6^n&1\le n\le8&8\\
2^a3^b&M<6^9,\ a\ne b,\ \gcd(a,b)<6&143\\
2^a5^b&M<9\,000\,000,\ (a,b)\ne(1,1),\ \gcd(a,b)<6&97\\
2^a7^b&M<9\,000\,000,\ (a,b)\ne(1,1),\ \gcd(a,b)<6&78\\
3^a5^b&M<2^{22},\ (a,b)\ne(1,1),\ \gcd(a,b)<6&54\\
B^6&B\in\{10,12,14\}&3
\end{array}
\]
These rows are not summed into a single disjoint inventory: their
roles in the dependency proof differ.
The first row plus the six-radical cutoff proves the base-six
all-length statement.  Every two-prime base at length at least six
is then covered: at length six, bases at least fifteen have modulus
greater than \(6^9\), base six is already covered, and the other
bases below fifteen are exactly \(10,12,14\), the last row.
For lengths at least seven, base ten reaches its even cutoff,
and bases at least twelve reach \(6^9\).
Thus a modulus with \(\gcd(a,b)\ge6\) is covered by reindexing at
that same modulus.  The middle four rows, the squarefree-semiprime
theorem, and their respective large-modulus cutoffs now exhaust
every valuation in radicals \(6,10,14,15\).

\subsection{Certificate arithmetic and exhaustive final partition}
\label{ap:finite:partition}

The finite certificates use disjoint original gcd strata.
For \(g\mid M\), \(g<M\), put \(m=M/g\) and
\[
 P_g(T)=\sum_{\substack{e\le T/g\\(e,D)=1}}\mu_{\mathrm{Mob}}(e)
 U_m(\lfloor T/(ge)\rfloor)U_D(\lfloor T/e\rfloor).
\]
Through \(\lfloor\sqrt M\rfloor\), the exact stratum CDF is
\(2P_g(T)-J_g(T)\), where ordered closed collision phases
\(Af+Ed=M\) are assigned to \(g=\gcd(A,M)\).
A centered inverse modulo \(m\), of magnitude
\(h\le H=\lfloor(m-1)/2\rfloor\), supplies numerator \(\pm g\).
If one prime \(v\) is missing from \(m\), repair a denominator
divisible by \(v\) to \(m-h\), reversing the numerator sign.
The repair remains a unit of the original \(M\).
For \(T\ge g\) its CDF is \(2U_D(\min(T,H))\), with the extra pure-child
term
\[
 2\max\{0,V_v(H)-V_v(\max(0,m-T-1))\},\qquad
 V_v(z)=\lfloor z/v\rfloor-\lfloor z/D\rfloor.
\]
Denominator one independently contributes
\(2U_m(\min(\lfloor T/g\rfloor,H))\).
At \(m=2\) each distribution is one atom at height \(g\), not a
doubled pair; both are zero for \(T<g\).
The circle-gap argument with \(p\) multiples of a state gives the
global cap \(M/p\): the index difference is less than \(p\), hence
full-unit, and its centered residue is at most \(M/p\).
The state \(M/p\) attains this cap.  Thus each stratum is completely
covered at
\[
 K_g=\min\{M/p,\ g\lfloor m/2\rfloor,\ \max(g,D_m)\}.
\]
\[
 D_m=\begin{cases}
 1,&m=2,\\
 H,&\operatorname{rad}(m)=D,\\
 \max(H,m-v),&m\hbox{ is a pure-prime power missing }v .
 \end{cases}
\]
The earliest radical-six stratum batch used the wider cap
\(\floor{m r_m/2}\), with \(r_m=1\) for a mixed child and \(r_m=v\)
for a pure child missing prime \(v\). This is also safe: in the mixed
case it is at least \(H\); in the pure case, \(v\ge2\) makes it at
least \(m\), hence at least every repaired denominator \(m-h\).
The singleton \(m=2\) remains one explicit atom at height \(g\).
Later batches use the refined \(D_m\) above. These safe recipes are
distinguished in the provenance map.

Take the maximum of these credits inside a stratum, sum only over
disjoint strata, add zero once, and combine by maximum with the
global far credit from \eqref{eq:finite:far-cover}.
On an ordered calendar containing all integers through \(\lfloor\sqrt M\rfloor\),
successive \(\lceil257T/256\rceil\) thereafter, all divisors/caps,
the first strict far threshold, and \(M/p\), cumulative maxima give
integer lower CDF values \(L_i\) ending at \(M\).
The histogram \(h_i=L_i-L_{i-1}\) dominates the true height even
between knots, by CDF monotonicity.  Therefore
\[
 \mu(M)\le
 \frac{2\sum_i h_i u(t_i)-M l(M)}{2M\,2^{64}}
 =\frac{Q_M}{2M\,2^{64}},
 \quad l(n)\le2^{64}\log n\le u(n).
\]
This doubled denominator is unchanged for odd \(M\).
The enclosures follow from binary range reduction and
\[
 0\le\log y-2\sum_{j=0}^{K-1}\frac{z^{2j+1}}{2j+1}
 \le\frac{2z^{2K+1}}{(2K+1)(1-z^2)},\qquad z=\frac{y-1}{y+1},
 \quad1\le y\le2 .
\]
The finite assertion is \(Q_M<0\) on every row of each declared set.
Large intermediate central counts rely on the exact formula and its
generator; a retained checker is not described as a complete second
implementation of every central count or of the true far survival.

We now use the previously proved squarefree-semiprime theorem and the
two-prime all-length theorem for lengths at least two.
If \(e=\gcd(a,b)\ge2\), set \(C=p^{a/e}q^{b/e}\), so \(M=C^e\);
this is precisely the same state space and witness set.
Together with the four complete radicals and \eqref{eq:finite:cutoffs},
the only conservative residual is
\[
 \begin{split}
 {\cal R}=\{(M,p,q,a,b):\ &\gcd(a,b)=1,\ (a,b)\ne(1,1),\\
 &pq\notin\{6,10,14,15\},\ M<C(p)\}.
 \end{split}
\]
\[
 C(p)=\begin{cases}
147456,&p=2,\\262144,&p=3,\\360000,&p\ge5 .
\end{cases}
\]
It contains all positive singleton and unequal valuations not paid by
another stated rule.  Unique factorization gives one ordered tuple
per physical modulus.  Since \(M\ge p^2q>p^3\), it is sufficient
to enumerate primes \(p\le71\), primes \(q<C(p)/p^2\), and positive
prime powers subject to the exact strict product test and these predicates.
A second enumeration by factorization gives the same full list.

Inside this residual the ratio rule \eqref{eq:finite:ratio} has only the
following possible supports:
\[
\begin{array}{c|c}
(p,a,b)&\hbox{prime }q\hbox{ range}\\ \hline
(2,2,1)&4900\le q<36864\\
(2,3,1)&2450\le q<18432\\
(2,4,1)&3600\le q<9216\\
(3,2,1)&11700\le q<262144/9\\
(5,2,1)&10000\le q<14400
\end{array}
\]
Indeed \(q\ge c(p)p^a\) and \(p^aq<C(p)\) imply
\(c(p)p^{2a}<C(p)\), while \(a=1,b=1\) is excluded.
The lower endpoints also impose \(M\ge19600\).
This proves that the rule excludes no unseen exponent orientation.

The exact coverage table is
\[
\begin{array}{l|r|l}
\hbox{stage}&\hbox{rows}&\hbox{mathematical reason}\\ \hline
\hbox{conservative residual }{\cal R}&19239&
 \hbox{complete original predicate}\\
\hbox{ratio-paid part}&7857&\eqref{eq:finite:ratio}\\
\hbox{after the ratio filter}&11382&19239-7857\\
\hbox{analytically certified part}&6972&
 \hbox{negative upper from the five complete bounds}\\
\hbox{finite remainder}&4410&11382-6972\\
\hbox{both primes at least five}&1087&\hbox{every retained }Q_M<0\\
\hbox{smaller prime three}&1422&\hbox{every retained }Q_M<0\\
\hbox{smaller prime two}&1901&\hbox{every retained }Q_M<0\\ \hline
\hbox{remaining two-prime obligations}&0&1087+1422+1901=4410
\end{array}
\]
The \(6972\) analytic rows are a conservative certificate partition,
not a claim that all other exact analytic bounds or true means are
nonnegative.  Each row retains all five bound enclosures.
For \(\alpha\), the rational tails
\[
 \sum_{j=1}^{128}j^{-2}+\frac1{129}<\zeta(2)
 <\sum_{j=1}^{128}j^{-2}+\frac1{128}
\]
are inverted outward.  Integer square-root floors give rational
\(x_L\le\sqrt M\le x_U\) with the same fixed \(k\) and \(t_U<1\).
The gains use the lower bounds
\[
 \delta_L=\min_{z\in\{t_L,t_U\}}\frac{1+2z-3z^2}{2}
                         +\frac{t_L^3}{x_U},\qquad
 G_L=\max(1,\,2x_L-M+k^2).
\]
The first uses concavity, the second the exact identity
\(G_x=1+(1-t)(2k-1+t)\ge1\).
Their denominator is the exact \(M\), never an approximate root squared.
After lowering \(\alpha\), write
\[
 \Delta_L=C_L+s_L\ell_\star,\quad
 C_L=\alpha_L(5\rho/8+3/8)+\rho/2-1/2-3u^2/4,\quad
 s_L=\alpha_L\rho-1/2+3u^2/4 ,
\]
where \(\ell_\star=\ell_L\) for \(s_L\ge0\) and \(\ell_U\) otherwise.
Positive error coefficients are enclosed upward, while their own
nonnegative gains are enclosed downward.  Strictly negative resulting
uppers pay the analytic rows.  No sign-selected domain or interpolation
is used.

Every original two-prime modulus is thus covered by a preliminary
theorem, an explicitly proved reduction, or one of the complete finite
certificate sets.  We conclude \(\mu(p^a q^b)<0\) for all distinct
primes and positive valuations.
For \(B=p^\gamma q^\delta\), \(\gamma,\delta\ge1\), \(N\ge1\),
put \(M=B^N\).  The original height, state measure and full-unit marks
are identical, and \(\log M=N\log B\).  Division by \(\log B>0\)
therefore gives
\[
 \mathbb E\log_B R_{(B,N)}<N/2 .
\]
This conclusion is restricted to exactly two prime divisors.
It neither transfers to bases with three prime factors nor asserts
monotonicity, optimal margins, exact true means or novelty.

\begin{thebibliography}{99}
\bibitem{binary}
O.~Babanskyy,
\emph{Rational complexity of binary strings: exact mean, finite
distribution and logarithmic variance}, version 0.2.0, 2026.
\url{https://github.com/aconsciousfractal/FCIG-Exact-Mean-Bound-for-the-Rational-Complexity-of-Binary-Strings}.
\bibitem{survey}
H.~Lin, L.~Xiao and Z.~Chen,
\emph{A brief view on the linear and 2-adic complexities of pseudorandom
binary sequences}, AIMS Mathematics 11(7) (2026), 20267--20290.
\href{https://doi.org/10.3934/math.2026823}{doi:10.3934/math.2026823}.
\bibitem{GK}
M.~Goresky and A.~Klapper,
\emph{Algebraic Shift Register Sequences}, Cambridge University Press, 2012.
\href{https://doi.org/10.1017/CBO9781139057448}{doi:10.1017/CBO9781139057448}.
Author draft (2009): \url{https://www.cs.uky.edu/~klapper/pdf/algebraic.pdf}.
\bibitem{CW}
Z.~Chen and A.~Winterhof,
\emph{Probabilistic results on the 2-adic complexity},
Designs, Codes and Cryptography 93 (2025), 2191--2203.
\href{https://doi.org/10.1007/s10623-025-01592-1}{doi:10.1007/s10623-025-01592-1};
preprint v1: \url{https://arxiv.org/abs/2501.16785}.
\bibitem{CostelloWatts}
F.~Costello and P.~Watts,
\emph{A computational upper bound on Jacobsthal's function},
preprint version 2, 2012.
\url{https://arxiv.org/abs/1208.5342v2}.
\bibitem{BHS}
F.~P.~Boca, B.~Heersink and P.~Spiegelhalter,
\emph{Gap distribution of Farey fractions under some divisibility constraints},
Integers 13 (2013), A44.
\url{https://arxiv.org/abs/1301.0277}.
\bibitem{Heersink}
B.~Heersink,
\emph{Equidistribution of Farey sequences on horospheres in covers of
\(\mathrm{SL}(n+1,\mathbb Z)\backslash\mathrm{SL}(n+1,\mathbb R)\)
and applications},
Ergodic Theory and Dynamical Systems 41 (2021), 471--493.
\url{https://arxiv.org/abs/1712.03258}.
\bibitem{Marklof}
J.~Marklof, \emph{Smallest denominators},
Bulletin of the London Mathematical Society 56 (2024), 1920--1938.
\href{https://doi.org/10.1112/blms.13034}{doi:10.1112/blms.13034};
\url{https://arxiv.org/abs/2310.11251}.
\bibitem{MarklofLog}
J.~Marklof, \emph{The log moments of smallest denominators},
Integers 24 (2024), A55.
\href{https://doi.org/10.5281/zenodo.12167470}{doi:10.5281/zenodo.12167470}.
\end{thebibliography}
\end{document}
